\documentclass[11pt]{article}

\usepackage[T1]{fontenc}
\usepackage[utf8]{inputenc}
\usepackage{lmodern}
\usepackage{microtype}
\usepackage[margin=0.82in]{geometry}
\usepackage{amsmath,amssymb,mathtools,bm}
\usepackage{booktabs,longtable,tabularx,array,multirow}
\usepackage{enumitem}
\usepackage{xcolor}
\usepackage{listings}
\usepackage{siunitx}
\usepackage[numbers,sort&compress]{natbib}
\usepackage{hyperref}
\usepackage[nameinlink,noabbrev]{cleveref}
\usepackage{fancyhdr}
\usepackage{lastpage}

\hypersetup{
  colorlinks=true,
  linkcolor=blue!55!black,
  citecolor=blue!55!black,
  urlcolor=blue!55!black,
  pdftitle={Physics and Implementation of the ss-RF Simulation Package},
  pdfsubject={Spin-1 Pake doublet, ss-RF, spectral spin diffusion, DNP, tensor-polarization design},
  pdfkeywords={spin-1, deuteron, Pake doublet, DNP, ss-RF, spin diffusion, tensor polarization, adjoint optimization}
}

\pagestyle{fancy}
\fancyhf{}
\lhead{ss-RF simulation package}
\rhead{Physics and Implementation Note}
\cfoot{\thepage\ of \pageref{LastPage}}
\setlength{\headheight}{14pt}

\setcounter{tocdepth}{2}
\setcounter{secnumdepth}{3}
\renewcommand{\arraystretch}{1.18}
\setlist[itemize]{topsep=3pt,itemsep=2pt,parsep=1pt}
\setlist[enumerate]{topsep=3pt,itemsep=2pt,parsep=1pt}

\definecolor{codebg}{RGB}{248,248,248}
\definecolor{codeframe}{RGB}{210,210,210}
\lstset{
  basicstyle=\ttfamily\small,
  backgroundcolor=\color{codebg},
  frame=single,
  rulecolor=\color{codeframe},
  breaklines=true,
  columns=fullflexible,
  showstringspaces=false
}

\newcommand{\Iplus}{I^{+}}
\newcommand{\Iminus}{I^{-}}
\newcommand{\Qb}{Q_{\mathrm B}}
\newcommand{\dd}{\mathrm{d}}
\newcommand{\e}{\mathrm{e}}
\newcommand{\Tr}{\operatorname{Tr}}
\newcommand{\clip}{\operatorname{clip}}
\newcommand{\avg}[1]{\left\langle #1\right\rangle}
\newcommand{\kb}{\bar{k}}
\newcommand{\code}[1]{\texttt{#1}}
\newcolumntype{L}[1]{>{\raggedright\arraybackslash}p{#1}}

\title{\textbf{ss-RF Simulation Package}\\[0.4em]
\large Spin-1 Population Physics, Spectral Spin Diffusion, Tensor-Polarization Design,\\
and Implementation of the Production Release}
\author{Technical documentation for the UVA simulation software}
\date{}

\begin{document}
\sloppy
\maketitle

\begin{abstract}
This note documents the physical model, its derivation, its numerical implementation, and the
tools built on the selective semi-saturation radio-frequency (ss-RF) simulation package for
spin-1 continuous-wave NMR, in its production release (\texttt{ssRF}). The modelled observable
is the overlapping Pake doublet of a quadrupolar spin-1 ensemble, with deuterons in solid
polarized targets (d-butanol, ND$_3$) as the motivating case. The state is resolved into the
three Zeeman-level populations of a set of orientation/frequency packets and, optionally, of
two chemical sites (C--D and O--D) together with the inverse temperature of the nuclear dipolar
reservoir.

The model contains:
\begin{itemize}
\item an analytic broadened Pake lineshape and a fitted two-site butanol lineshape;
\item signed vector and tensor polarization with exact direct-hole and mirror-peak bookkeeping;
\item an RF operator with a finite homogeneous linewidth, driven by arbitrary per-bin, time-scheduled
programs, with separately tunable per-spin coupling and width profiles;
\item spectral spin diffusion as mass-action pair currents for the three spin-1 flip-flop channels
(same-branch, cross-branch, and an effective double-quantum channel) on a nonlocal Lorentzian
zero-quantum kernel;
\item DNP and spin-lattice relaxation;
\item instantaneous adiabatic fast passage;
\item a first-order Provotorov coupling to the dipolar reservoir.
\end{itemize}
Conservation of spins, of vector polarization and (where appropriate) of energy follows from the
structure of the currents rather than from projections. The unique stationary state of the
internal dynamics, a Boltzmann line at the remaining polarization, is derived rather than imposed.

The note then documents:
\begin{itemize}
\item the time integrator;
\item the tensor-polarization optimizer, a discrete-adjoint, constrained search for the RF program
that maximizes $Q$ under power and energy limits;
\item the calibrated d-butanol materials;
\item a command-line training-data generator, including a closed-loop mode that reproduces the
operating cycle;
\item the machine-learning models trained on its output;
\item the graphical interface, the regression tests, and a staged calibration procedure.
\end{itemize}
It also quantifies the peak--pedestal saturation balance with the calibrated materials. It
supersedes the first-release note, \emph{Real-Time ss-RF Simulation Package:
Spin-Capacity-Weighted DNP, Depolarization, and Recovery}, and \cref{sec:changes} lists what
changed and why. Exact spin-1 bookkeeping is kept separate from phenomenological closures
throughout.
\end{abstract}

\tableofcontents
\clearpage

%==================================================================================================
\section{Scope, model status, and central physical picture}
\label{sec:scope}

The ss-RF simulation package is a forward model for the evolution of a spin-1 Pake-doublet NMR
signal under selective RF, adiabatic fast passage (AFP), DNP, spectral spin diffusion, and
spin-lattice relaxation. It also contains an inverse model: a designer that computes the RF
program which maximizes the tensor polarization from a given state. It is not a first-principles
many-body calculation. It is a physically constrained, experimentally parameterizable simulator,
calibrated for deuterated butanol, intended to answer questions such as:

\begin{itemize}
  \item What happens to the two overlapping absorption components when RF of finite linewidth is applied at one or many physical frequencies?
  \item How large is the direct hole in each branch, where do the corresponding mirror peaks appear, and how does a finite RF linewidth modify the classical one-half mirror rule?
  \item How does a hole recover with the microwaves off, when vector polarization removed by ss-RF cannot be recreated by spin diffusion, and how does recovery depend on which flip-flop channels operate?
  \item How does the answer change when DNP is on and supplies polarization from the electron-spin reservoir?
  \item Which RF program, within given power and energy limits, maximizes the tensor polarization from the present state, and how does that program change as the state evolves?
  \item Why does a hole near a Pake peak need a different RF power than one in the pedestal, and which material parameters decide this?
  \item Which rates must be measured, as functions of frequency, polarization, RF power, temperature and material, to make the predictions quantitative?
\end{itemize}

The physical basis follows the spin-1 population and Pake-doublet formalism developed for
RF-manipulated solid targets in Refs.~\cite{ClementKeller2023,Keller2026,KellerCrabbDay2020,Keller2017},
together with standard solid-state NMR, spin-temperature and DNP theory
\cite{Abragam1961,Slichter1990,Wasylishen2012,Goldman1970,AbragamGoldman1978,AbragamGoldman1982}.
The characteristic powder-pattern density originates in the orientation-to-frequency mapping
first associated with Pake-type spectra \cite{Pake1948}.

The spin-diffusion architecture follows the orientation-resolved model introduced in
\emph{Selective semi-saturation measurement of spin-1} \cite{Keller2026}. That work formulates
the level populations as \(n_i(\theta,t)\), lets neighbouring orientation packets exchange
polarization, and proposes a \(T_2\)-limited spectral-overlap kernel. The first release reduced
this to a local relaxation plus a nearest-neighbour flux. The production release implements the
spectral-overlap idea directly. Every pair of packets exchanges population through the
microscopic spin-1 flip-flop channels, at a rate set by their zero-quantum spectral overlap and
by the instantaneous availability of partners in the right levels \cite{SuterErnst1985,Khutsishvili1966,Eckman1982}.

The package rests on six central statements.

\begin{enumerate}
  \item A spin-1 packet has three populations \((n_{+},n_{0},n_{-})\).  The two allowed single-quantum absorption intensities are population differences, \(\Iplus\propto n_{+}-n_{0}\) and \(\Iminus\propto n_{0}-n_{-}\).
  \item At one physical frequency \(R\) in the overlap region, \(\Iplus(R)\) and \(\Iminus(R)\) generally come from two different orientation packets.  A common RF tone therefore drives two distinct transitions with one common applied power.
  \item Each orientation packet is observed twice: once in one absorption branch at \(R\), and once in the other branch at \(-R\).  The mirror peak is therefore not an independent reservoir; it is another observable projection of the same three populations.
  \item ss-RF equalizes the addressed pair of levels.  It depolarizes the usual monotone spin-temperature populations; for a general non-Boltzmann state its effect on \(P\) and \(Q\) follows from the population currents.  Internal spin diffusion redistributes polarization but cannot create vector polarization.  DNP and spin-lattice relaxation are the external sources and sinks.
  \item An equal-width frequency bin does not contain the same number of spins at every \(R\).  This \emph{capacity} \(\mu_k\), set by the Pake density, enters every current through the populations.  It is distinct from the \emph{rate per spin}.  In the production release any additional dependence of a per-spin rate on frequency, orientation or site is an explicit, separately adjustable mechanism profile.  It is never a hidden second use of the density.
  \item Recovery is a set of pair processes with definite matrix elements and definite conservation laws.  Same-branch flip-flops conserve both \(P\) and \(Q\); cross-branch flip-flops conserve \(P\) but change \(Q\); double-quantum exchange moves vector polarization between packets without frequency matching.  Which channels dominate is the largest open material question for tensor enhancement.
\end{enumerate}

\begin{quote}
\textbf{Scientific status of the simulation closure.}  The following are structurally fixed by spin-1 quantum mechanics:
\begin{itemize}
  \item the packet capacity \(\mu_k\), the population bookkeeping and the mirror relation;
  \item the definitions of \(P\) and \(Q\);
  \item the equality of the two single-quantum RF matrix elements;
  \item the channel structure and conservation laws of the flip-flop currents;
  \item the uniqueness of the Boltzmann stationary state.
\end{itemize}
The following are phenomenological closures with calibrated or adjustable parameters:
\begin{itemize}
  \item the zero-quantum overlap shape and width, and the effective double-quantum ratio;
  \item the per-spin RF selectivity;
  \item the DNP target law;
  \item the first-order dipolar-reservoir coupling.
\end{itemize}
For d-butanol the lineshape, DNP build-up, RF strength and exchange were fitted to UVA target data.  Per-spin RF selectivity and the cross-branch ratio are not constrained by present data (\cref{sec:materials}).
\end{quote}

%==================================================================================================
\section{Spin-1 energy levels and the quadrupolar orientation dependence}
\label{sec:levels}

\subsection{High-field Hamiltonian}

A deuteron has spin \(I=1\), so in a strong holding field it has three Zeeman sublevels labelled
\(m=+1,0,-1\). In a non-cubic solid, the nuclear quadrupole moment couples to the local
electric-field-gradient (EFG) tensor. To first order in the quadrupolar coupling,
\begin{equation}
  \hat H
  = -\hbar\omega_d \hat I_z
  + \hbar\omega_q\,
  \Lambda(\theta,\phi)
  \left(3\hat I_z^2-I(I+1)\right),
  \label{eq:hamiltonian}
\end{equation}
where \(\omega_d\) is the deuteron Larmor angular frequency, \(\omega_q\) sets the quadrupolar
scale, and
\begin{equation}
  \Lambda(\theta,\phi)
  = 3\cos^2\theta-1+\eta\sin^2\theta\cos 2\phi.
  \label{eq:lambda}
\end{equation}
The angle \(\theta\) is measured between the EFG principal axis and the holding field, and
\(\eta\) is the EFG asymmetry \cite{Abragam1961,CohenReif1957,Wasylishen2012}. For
\(\hat H_Q\propto e^2qQ\), \(\omega_q=e^2qQ/(8\hbar)\) in the normalization of
\cref{eq:hamiltonian}. Equivalent normalizations occur in the literature; the simulation is
built around the dimensionless frequency coordinate below.

For \(I=1\), \(3\hat I_z^2-I(I+1)\) takes the value \(+1\) for \(m=\pm1\) and \(-2\) for
\(m=0\). The level energies are therefore
\begin{equation}
  E_{\pm1}=\mp\hbar\omega_d+\hbar\omega_q\Lambda,\qquad E_0=-2\hbar\omega_q\Lambda,
\end{equation}
and the two allowed magnetic-dipole transitions,
\begin{equation}
  m=+1\leftrightarrow 0,
  \qquad
  m=0\leftrightarrow -1,
\end{equation}
occur at \(\omega_d\mp3\omega_q\Lambda\). The \(\Delta m=2\) transition \(+1\leftrightarrow-1\)
lies at \(2\omega_d\) with \emph{no} first-order quadrupolar shift. This fact underlies the
double-quantum exchange channel of \cref{sec:dq}. Because the quadrupolar contribution depends
on \(\theta\), the transition frequencies vary across a polycrystalline or amorphous sample.
Powder averaging therefore produces two broad, overlapping absorption components rather than two
sharp lines \cite{ClementKeller2023,Keller2026}.

\paragraph{Rotating-frame level energies used by the code.}
Remove the Larmor term in the rotating frame, subtract a constant, and measure energy in units
of \(3\hbar\omega_q\). The level energies of a packet whose \(\Iplus\) transition lies at \(R\)
(\cref{eq:Rplus}) are then
\begin{equation}
  E_{+}=E_{-}=-R,\qquad E_0=0 .
  \label{eq:levels-rot}
\end{equation}
The \(\Iplus\) transition is at \(E_0-E_+=R\), the \(\Iminus\) transition at \(E_--E_0=-R\),
and the double-quantum transition at \(E_--E_+=0\). The total offset (non-Larmor) energy of the
spin system is
\begin{equation}
  E_{\rm spin}=-\sum_k R_k\,(n_{k+}+n_{k-}),
  \label{eq:Espin}
\end{equation}
which the extended engine balances against the dipolar reservoir (\cref{sec:extended}).

\paragraph{Magnetic-dipole matrix elements.}
In the high-field limit a transverse RF field couples through \(\hat I_x\), with
\begin{equation}
  |\langle0|\hat I_x|+1\rangle|^2
  =|\langle-1|\hat I_x|0\rangle|^2=\frac12,
  \qquad
  |\langle0|\hat I_+|-1\rangle|^2=|\langle+1|\hat I_+|0\rangle|^2=2 .
  \label{eq:rf-matrix-elements}
\end{equation}
Both single-quantum transitions therefore have the same coupling, and the EFG angle introduces no
leading-order suppression of the magnetic-dipole drive at the powder peak \cite{Abragam1961}.

\subsection{Dimensionless frequency coordinate and software branch convention}

The frequency coordinate is
\begin{equation}
  R = \frac{\omega-\omega_d}{3\omega_q}.
  \label{eq:Rdef}
\end{equation}
The simulation software adopts the branch mapping of Ref.~\cite{ClementKeller2023}:
\begin{align}
  R_{+}(\theta) &= 1-3\cos^2\theta,
  & -2\le R_{+}\le 1,
  \label{eq:Rplus}\\
  R_{-}(\theta) &= 3\cos^2\theta-1=-R_{+}(\theta),
  & -1\le R_{-}\le 2,
  \label{eq:Rminus}
\end{align}
for \(\eta=0\); with asymmetry,
\(R_+=1-3\cos^2\theta-\eta\sin^2\theta\cos2\phi\). The labels \(+\) and \(-\) identify the two
absorption branches as used by the code. The main consequences are independent of sign
convention:

\begin{itemize}
  \item \(\theta=90^\circ\) maps to the \(\Iplus\) horn at \(R=+1\) and to the \(\Iminus\) horn at \(R=-1\).
  \item \(\theta=0\) maps to \(R=-2\) for \(\Iplus\) and \(R=+2\) for \(\Iminus\).
  \item The ``magic'' orientation \(\cos^2\theta=1/3\) maps to \(R=0\) in both branches.
  \item One orientation packet has two mirror-related spectral appearances, at \(R\) and \(-R\).
\end{itemize}

In physical units, \(R=1\) is \(3\omega_q/2\pi=\tfrac38e^2qQ/h\). That is
\(\approx\SI{63}{kHz}\) for C--D bonds with \(e^2qQ/h\approx\SI{167}{kHz}\). The O--D site of
butanol has a larger coupling, so its line extends to \(|R|\approx2.7\) on the C--D scale. The
broadened lineshape has tails beyond the strict first-order domains, so the grid extends beyond
\(\pm2\): \(\pm3\) for the single-site demonstration material and \(\pm6.19\) for the d-butanol
materials, which must contain the O--D component and its tails.

\subsection{Why a Pake horn contains more spin capacity}

For an isotropic powder, the number of EFG axes in a polar-angle interval is proportional to the
solid-angle measure
\begin{equation}
  \dd N \propto \sin\theta\,\dd\theta.
\end{equation}
Equivalently, \(u=\cos\theta\) is uniformly distributed on the hemisphere. For the \(\Iplus\)
branch,
\begin{equation}
  R_{+}=1-3u^2,
  \qquad
  u=\sqrt{\frac{1-R_{+}}{3}}.
\end{equation}
Therefore the ideal, unbroadened density per unit \(R\) is
\begin{equation}
  \frac{\dd N}{\dd R_{+}}
  \propto
  \left|\frac{\dd u}{\dd R_{+}}\right|
  = \frac{1}{2\sqrt{3}}\frac{1}{\sqrt{1-R_{+}}},
  \qquad
  \int_{-2}^{1}\frac{\dd R_+}{2\sqrt3\sqrt{1-R_+}}=1 .
  \label{eq:pakejacobian}
\end{equation}
This Jacobian diverges at \(R_{+}=1\), corresponding to \(\theta=90^\circ\). Broadening makes the
horn finite, but the physical point remains: many orientations are compressed into a small
frequency interval. The \(\Iminus\) branch is the mirror image and has its horn at \(R=-1\).
This powder-density enhancement is the geometric origin of the Pake horns
\cite{Pake1948,Abragam1961}.

The Jacobian counts \emph{spins per unit frequency}. It does not set a \emph{rate per spin}. The
golden-rule rate of \cref{sec:rf-golden} contains no factor of the density, and a flip-flop
current between two packets contains the capacities of both packets exactly once, through their
populations. Any further dependence of a per-spin rate on frequency must come from a separate
physical mechanism, for example an orientation-dependent \(T_2\). The production release models
such mechanisms with the explicit profiles of \cref{sec:profiles}.

%==================================================================================================
\section{Populations, vector and tensor polarization, and NMR observables}
\label{sec:observables}

\subsection{Three-level populations}

Let the normalized total populations of the three Zeeman levels be
\begin{equation}
  n_{+}+n_{0}+n_{-}=1.
\end{equation}
The vector and tensor polarizations are
\begin{align}
  P &= n_{+}-n_{-},
  \label{eq:Pdef}\\
  Q &= n_{+}-2n_{0}+n_{-}=1-3n_{0}.
  \label{eq:Qdef}
\end{align}
Solving for the level fractions gives
\begin{align}
  p_{+}(P,Q) &= \frac13+\frac{P}{2}+\frac{Q}{6},
  \label{eq:pplus}\\
  p_{0}(P,Q) &= \frac{1-Q}{3},
  \label{eq:pzero}\\
  p_{-}(P,Q) &= \frac13-\frac{P}{2}+\frac{Q}{6}.
  \label{eq:pminus}
\end{align}
The physical region is the triangle in \((P,Q)\) for which all three fractions are nonnegative:
\(Q\le1\), \(Q\ge 3|P|-2\). The backend accepts an optional initial \(Q_0\); materials and the
GUI initialize with the Boltzmann value unless a non-Boltzmann state is loaded from a snapshot.

\subsection{Boltzmann or single-spin-temperature relation}

For a spin-1 system with one Zeeman spin temperature, \(p_m\propto\e^{\beta m}\). Writing
\(r=\e^{\beta}\),
\begin{equation}
  P=\frac{r-r^{-1}}{r+1+r^{-1}},\qquad
  Q=\frac{r+r^{-1}-2}{r+1+r^{-1}} .
\end{equation}
Eliminating \(r\) (with \(x=r+r^{-1}\), so \(Q=(x-2)/(x+1)\) and \(P^2=(x^2-4)/(x+1)^2\)) gives
\begin{equation}
  Q_{\mathrm B}(P)
  = 2-\sqrt{4-3P^2}.
  \label{eq:QB}
\end{equation}
For example, \(Q_{\rm B}(0.50)=0.1972\). This relation is used to construct the default initial
state, the DNP and \(T_1\) targets, and the targets in the data generator. It is \emph{not}
imposed on the evolving state, and it is no longer used as a recovery reference. In the
production release, the approach of an isolated state to the Boltzmann line follows from the
flip-flop currents (\cref{sec:stationary}). During RF manipulation \(Q(t)\) can differ from
\(Q_{\rm B}[P(t)]\), and the GUI displays both.

The branch population differences in a uniform Boltzmann state are
\begin{align}
  p_{+}-p_{0} &= \frac{P+Q}{2},
  \label{eq:branchdiffplus}\\
  p_{0}-p_{-} &= \frac{P-Q}{2}.
  \label{eq:branchdiffminus}
\end{align}
Their sum is \(P\) and their difference is \(Q\). The corresponding equilibrium branch-area ratio
is
\begin{equation}
  r(P)=\frac{I^{+}}{I^{-}}
  =\frac{\sqrt{4-3P^2}+P}{2-2P},
  \label{eq:branchratio}
\end{equation}
which is the ratio used in the 2023 formalism \cite{ClementKeller2023}.

\subsection{Absorption intensities}

For an orientation packet at \(\theta\), the absorptive components are
\begin{align}
  \Iplus(\theta)&=C\,[\rho_{+}(\theta)-\rho_{0}(\theta)],
  \label{eq:IplusTheta}\\
  \Iminus(\theta)&=C\,[\rho_{0}(\theta)-\rho_{-}(\theta)],
  \label{eq:IminusTheta}
\end{align}
where \(C\) is the NMR calibration factor and the \(\rho_m\) are population densities.
Integrated over the full branches,
\begin{align}
  P &= C\,(I^{+}+I^{-}),
  \label{eq:parea}\\
  Q &= C\,(I^{+}-I^{-}).
  \label{eq:qarea}
\end{align}
These area relations are definitions and hold for any state. What fails out of equilibrium is
the inference of \(Q\) from \(P\) using \cref{eq:QB} \cite{ClementKeller2023,Keller2026}. That
failure is why a learned polarization model is useful for manipulated states (\cref{sec:ml}).

Signed polarization is handled naturally. When \(P<0\), the ordering of the three populations
reverses and both absorptive population differences can become negative. The software plots
signed branch intensities, and every rate equation is written for signed populations.

%==================================================================================================
\section{Packet space, physical-frequency overlap, and mirror bookkeeping}
\label{sec:packets}

\subsection{State coordinate used by the simulation software}

The dynamic state is stored on a one-dimensional packet coordinate
\begin{equation}
  x_k \equiv R_{+}(\theta_k),
  \qquad k=0,\ldots,N-1,
\end{equation}
on a uniform grid of width \(\Delta R\). Each packet contains a three-component population vector
\begin{equation}
  \bm n_k=(n_{k+},n_{k0},n_{k-})^{\mathsf T},
  \qquad
  n_{k+}+n_{k0}+n_{k-}=\mu_k,
  \qquad
  \sum_k\mu_k=1.
  \label{eq:packetstate}
\end{equation}
The packet is observed as
\begin{align}
  \Iplus\text{ at physical }R=x_k,
  \qquad
  \Iminus\text{ at physical }R=-x_k.
  \label{eq:packetmirror}
\end{align}
Thus the two branch appearances are not separate states.

\paragraph{Symmetric grid and the mirror index.}
The production release requires the grid to be symmetric about zero, to \(10^{-12}\). Then the
packet whose \(\Iplus\) line is at \(-x_k\) is
\begin{equation}
  \kb\equiv N-1-k,\qquad x_{\kb}=-x_k,
  \label{eq:mirrorindex}
\end{equation}
and every mirror lookup is an exact index reversal, with no interpolation.
\code{IdealBinModel.reset} enforces the symmetry. The d-butanol materials use \(N=500\) on
\([-6.186,6.186]\) (\(\Delta R=0.0248\)); the demonstration material uses \(N=500\) on
\([-3,3]\).

\subsection{Two different packets at one physical RF frequency}

At a physical frequency \(R_b\) the two overlapping components are
\begin{align}
  \Iplus(R_b) &: \text{packet }k_{+}\text{ with }x_{k_{+}}\simeq R_b,
  \label{eq:kplus}\\
  \Iminus(R_b) &: \text{packet }k_{-}=\overline{k_+}\text{ with }x_{k_{-}}\simeq -R_b.
  \label{eq:kminus}
\end{align}
These are generally two different values of \(\theta\). A single RF tone has one amplitude, but
it acts on both transitions because both are resonant at \(R_b\). With the matrix elements of
\cref{eq:rf-matrix-elements}, the two transitions have the same golden-rule rate per spin
(\cref{sec:rf-golden}). Their currents differ because their population differences differ, and
because of any explicit per-spin profile.

The mirror observables are
\begin{align}
  \Iminus(-R_b)&\text{ from packet }k_{+},
  \label{eq:mirrorminus}\\
  \Iplus(-R_b)&\text{ from packet }k_{-}.
  \label{eq:mirrorplus}
\end{align}
This bookkeeping is essential: direct holes and mirror peaks must be calculated from shared
populations, not evolved as four independent signals.

%==================================================================================================
\section{Static lineshape and packet capacity}
\label{sec:lineshape}

\subsection{Analytic broadened single-site branch function}

The single-site lineshape (\code{line\_model = single\_pake}) is the analytic broadened Pake
function of Refs.~\cite{Dulya1997,ClementKeller2023}. Define
\begin{align}
  s &\equiv \eta\cos 2\phi,\\
  X_2(R,\epsilon) &\equiv \sqrt{\Gamma_{\rm line}^2+(1-\epsilon R-s)^2},\\
  X(R,\epsilon) &\equiv \sqrt{X_2(R,\epsilon)},\\
  Y &\equiv \sqrt{3-s},\\
  \cos\alpha &\equiv \frac{1-\epsilon R-s}{X_2(R,\epsilon)}.
\end{align}
Then the branch density is
\begin{align}
 f(R,\epsilon)
 =\frac{1}{2\pi X}
 \Bigg\{&2\cos\frac{\alpha}{2}
 \left[
 \tan^{-1}\left(
 \frac{Y^2-X_2}{2YX\sin(\alpha/2)}
 \right)+\frac{\pi}{2}
 \right]
 \nonumber\\
 &+\sin\frac{\alpha}{2}
 \ln\left[
 \frac{Y^2+X_2+2YX\cos(\alpha/2)}
      {Y^2+X_2-2YX\cos(\alpha/2)}
 \right]\Bigg\}.
 \label{eq:analyticpake}
\end{align}
The branch \(\epsilon=+1\) has its sharp horn on the positive-\(R\) side, and
\(f(R,-1)=f(-R,+1)\). The demonstration material uses \(\Gamma_{\rm line}=0.05\) and
\(s=0.04\), approximate ND\(_3\) values \cite{ClementKeller2023}.

\subsection{Azimuth-averaged resolvent form for general asymmetry}
\label{sec:resolvent}

For the fitted materials (\code{line\_model = multisite\_butanol}), each site uses the full
azimuthal average. Write \(c=\cos\theta\) and \(y_\phi^2=3-\eta\cos2\phi\). Then the static
offset is \(x=1-\eta\cos2\phi-y_\phi^2c^2\). A Lorentzian broadening of half-width \(\gamma\) is
the imaginary part of the resolvent:
\begin{equation}
g_\gamma(x) = -\frac1\pi\,\mathrm{Im}\,\Bigl\langle\int_0^1\frac{\dd c}{1-\epsilon x-\eta\cos2\phi-y_\phi^2c^2+i\gamma}\Bigr\rangle_\phi
= -\frac1\pi\,\mathrm{Im}\,\Bigl\langle\frac{\operatorname{artanh}(y_\phi/r_\phi)}{y_\phi\,r_\phi}\Bigr\rangle_\phi,
\label{eq:resolvent}
\end{equation}
with \(r_\phi=\sqrt{1-\epsilon x-\eta\cos2\phi+i\gamma}\). The \(c\)-integral is elementary
(\(\int_0^1\dd c/(r^2-y^2c^2)=\operatorname{artanh}(y/r)/(yr)\)). Only the average over
\(\phi\in[0,\pi/2]\) is done numerically, by Gauss--Legendre quadrature
(\code{line\_phi\_points}; code default 64, 128 in the d-butanol materials). For \(\epsilon=-1\) the code conjugates the kernel, which
mirrors the branch. The form is exact for a Lorentzian convolution of the first-order powder
pattern, and has unit integrated absorption on the infinite axis.

\subsection{Two-site butanol lineshape}
\label{sec:twosite}

Butanol-\(d_{10}\) contains C--D and O--D deuterons with different quadrupole couplings. Let
\(a_{\rm CD}\) be the C--D scale (one unit of \(R\)), \(a_{\rm OD}=\rho_{\rm OD}a_{\rm CD}\) and
\(f_{\rm OD}\) the O--D fraction. The capacity density is
\begin{equation}
g(R) = \frac{1-f_{\rm OD}}{a_{\rm CD}}\,g_{\gamma/a_{\rm CD}}^{(\eta_{\rm CD})}\!\Bigl(\frac{R}{a_{\rm CD}}\Bigr)
+\frac{f_{\rm OD}}{a_{\rm OD}}\,g_{\gamma/a_{\rm OD}}^{(\eta_{\rm OD})}\!\Bigl(\frac{R}{a_{\rm OD}}\Bigr),
\label{eq:multisite}
\end{equation}
so that \(f_{\rm OD}\) is an integrated site weight and both sites share one physical width
\(\gamma\). The fit to the 2014 UVA spectra (held-out RMS 0.79\%) gives
\(\rho_{\rm OD}=1.355\), \(f_{\rm OD}=0.021\), \(\eta_{\rm CD}=0.065\), \(\eta_{\rm OD}=0.033\) and
\(\gamma=0.0236\).

Mechanisms that depend on orientation need each packet's angle and site composition. The code
inverts the \(\eta=0\) map per site,
\(\cos^2\theta_s(R)=\clip\bigl((1-R/a_s)/3,0,1\bigr)\), and records the O--D share
\(w_{\rm OD}(R)=g_{\rm OD}/(g_{\rm CD}+g_{\rm OD})\) (\code{packet\_site\_geometry}).

\subsection{Normalization and packet capacity}

The \(\epsilon=+1\) density is normalized on the simulation grid and converted to packet
fractions,
\begin{equation}
  \rho_k=\frac{g(x_k)}{\int g(x)\,\dd x},
  \qquad
  \mu_k = \frac{\rho_k\Delta R}{\sum_j\rho_j\Delta R}.
  \label{eq:muk}
\end{equation}
The initial packet populations are
\begin{equation}
  n_{km}(0)=\mu_k p_m(P_0,Q_0).
  \label{eq:initpacket}
\end{equation}
Thus \(\mu_k\) already changes the absolute number of spins and the absolute signal in each bin.
A bin with larger \(\mu_k\) contains more population to be moved by every process. The displayed
intensities are
\begin{align}
  \Iplus(x_k)&=\mathcal C_{\rm disp}\frac{n_{k+}-n_{k0}}{\Delta R},
  \label{eq:displayIplus}\\
  \Iminus(-x_k)&=\mathcal C_{\rm disp}\frac{n_{k0}-n_{k-}}{\Delta R}.
  \label{eq:displayIminus}
\end{align}
The display calibration \(\mathcal C_{\rm disp}\) affects only plotted units.

\subsection{Capacity is not identical to a microscopic rate}
\label{sec:capacity-vs-rate}

A central interpretive distinction is
\begin{equation}
  \underbrace{\text{observable current}}_{\dd I/\dd t}
  \sim
  \underbrace{\text{number of spins}}_{\mu_k}
  \times
  \underbrace{\text{population imbalance per spin}}_{p_a-p_b}
  \times
  \underbrace{\text{transition rate per spin}}_{W_k}.
  \label{eq:capacityratefactorization}
\end{equation}
Because \(n_{ka}=\mu_k p_{ka}\), the capacity is already present in the state. With a uniform
per-spin rate, a horn bin still produces a larger \emph{absolute} source or sink, because it
contains more spins. Its \emph{fractional} saturation proceeds at the same rate as anywhere else.

The first release multiplied every rate by an extra weight \(w_k\propto\rho_k^\alpha\). With
\(\alpha=1\) this made absolute horn currents scale roughly as \(\rho^2\). The production
release removes that double use:
\begin{itemize}
  \item RF and spin-diffusion currents contain the capacity once.
  \item Any further frequency, orientation or site dependence of a per-spin rate is an explicit
  mechanism profile with its own parameters (\cref{sec:profiles}).
  \item The old weight survives only as the DNP build-rate weighting (\code{capacity\_rate\_power},
  \cref{sec:dnp}) and in the legacy RF modes kept for comparison.
\end{itemize}

\begin{quote}
\textbf{Practical interpretation.}  A fractional quantity (hole depth relative to the local line,
fractional build-up, required RF power for a given fractional saturation) is independent of
\(\mu_k\) unless a per-spin rate depends on \(k\).  A difference between peak and pedestal in
such a quantity therefore measures a per-spin difference in RF response, in replenishment, or
in both.  \Cref{sec:peak-pedestal} shows how to separate them, and what the calibrated
materials predict.
\end{quote}

%==================================================================================================
\section{Complete rate model implemented in the simulation software}
\label{sec:ratemodel}

\subsection{Additive mechanism decomposition}

Without the extended engine (\cref{sec:extended}), the model evaluates the population derivative as
\begin{equation}
  \dot{\bm n}
  =\dot{\bm n}_{\rm RF}
  +\dot{\bm n}_{\rm SQ}
  +\dot{\bm n}_{\rm X}
  +\dot{\bm n}_{\rm DQ}
  +\dot{\bm n}_{\rm DNP}
  +\dot{\bm n}_{T_1},
  \label{eq:totalrate}
\end{equation}
with AFP applied as a discrete map between integration steps (\cref{sec:afp}). The terms have
distinct physical roles:

\begin{center}
\begin{tabularx}{0.96\textwidth}{>{\bfseries}L{0.20\textwidth} X L{0.24\textwidth}}
\toprule
Term & Physical role & Conserves\\
\midrule
ss-RF & Incoherent saturation of every transition within the homogeneous linewidth of each carrier & \(\mu_k\); not \(P\), not \(Q\)\\
Same-branch (SQ) & Flip-flops \(|{+}0\rangle\leftrightarrow|0{+}\rangle\) and \(|0{-}\rangle\leftrightarrow|{-}0\rangle\) between packets with overlapping transitions & \(\mu_k\), \(N^{\pm,0}\), \(P\), \(Q\)\\
Cross-branch (X) & Flip-flops \(|00\rangle\leftrightarrow|{+}{-}\rangle\) between packets whose \(\Iplus\) and \(\Iminus\) lines overlap & \(\mu_k\), \(P\); changes \(Q\)\\
Double quantum (DQ) & Effective \(|{+}{-}\rangle\leftrightarrow|{-}{+}\rangle\) exchange with no first-order frequency condition & \(\mu_k\), \(N^{\pm,0}\), \(P\), \(Q\)\\
DNP & Relaxation of every packet toward the Boltzmann fractions at \(P_{\rm sat}\) & \(\mu_k\)\\
\(T_1\) & Relaxation toward the Boltzmann fractions at \(P_{\rm eq}\) & \(\mu_k\)\\
\bottomrule
\end{tabularx}
\end{center}

Here \(N^m=\sum_kn_{km}\). No term resets the state toward a reference line, and no term is
projected after the fact. Every conservation property in the table follows from the form of the
currents.

\subsection{RF saturation: the golden-rule rate}
\label{sec:rf-golden}

A linearly polarized field \(2B_1\cos\omega t\) has a co-rotating component
\(\hat H_1=\tfrac{\hbar\omega_1}{2}(\hat I_+e^{-i\omega t}+\mathrm{h.c.})\) with
\(\omega_1=\gamma_dB_1\). Suppose the coherence between the levels \(a\) and \(b\) of a
transition decays exponentially at rate \(\Gamma\) (the homogeneous half-width, a \(T_2\)), and
the drive is weak, \(\omega_1\ll\Gamma\). The coherence then follows the populations
adiabatically, and eliminating it gives a rate equation
\(\dot n_a=-W(n_a-n_b)\) \cite{Bloch1946,Redfield1955,Abragam1961,Slichter1990} with
\begin{equation}
W = 2\pi\,\frac{\omega_1^2}{4}\,|\langle b|\hat I_+|a\rangle|^2\,g_{\rm L}(\Delta)
= \omega_1^2\,\frac{\Gamma}{\Gamma^2+\Delta^2},
\label{eq:Wgolden}
\end{equation}
where \(g_{\rm L}\) is the unit-area Lorentzian and the last step uses
\cref{eq:rf-matrix-elements}. Both spin-1 transitions have this same rate, twice the
spin-\(\tfrac12\) value. Equivalently, the population difference decays at
\(\lambda_{\rm RF}=2W\), which is the Bloch-rate form
\(\Omega_1^2T_2/(1+\Delta\omega^2T_2^2)\) with \(\Omega_1^2=2\omega_1^2\).

\paragraph{Validity.} In the code, carrier \(j\) has power \(U_j\propto B_1^2\), and the material
constant \(\kappa\) (\code{rf\_spectral\_strength}) absorbs the conversion
\(\omega_1^2\leftrightarrow U\) and the \(R\) unit:
\begin{equation}
W_{kj} = \kappa\,\frac{\Gamma}{\Gamma^2+(x_k-x_j)^2}\,U_j .
\label{eq:Wkj}
\end{equation}
For d-butanol, \(\kappa=0.0196\) and \(\Gamma=0.015\) give an on-resonance rate
\(\kappa/\Gamma=1.31\)\,s\(^{-1}\) per unit \(U\). With \(\Gamma\approx\SI{0.94}{kHz}\) in
physical units, this corresponds to \(\omega_1/2\pi\approx\SI{14}{Hz}\sqrt U\). That is orders of
magnitude below \(\Gamma\), so the weak-drive closure is self-consistent for every command the
designer produces. Coherent phenomena (nutation, spin locking, rotary saturation beyond the
Provotorov terms) are absent.

\subsection{The discrete RF operator}
\label{sec:rf-operator}

Packet \(k\) sees the RF at its \(\Iplus\) frequency \(x_k\) and at its \(\Iminus\) frequency \(-x_k\):
\begin{equation}
W^+_k = c_k\sum_j K_k(x_k-x_j)\,U_j,\qquad W^-_k = c_k\sum_j K_k(-x_k-x_j)\,U_j,\qquad
K_k(\delta)=\frac{\kappa\,\Gamma_k}{\Gamma_k^2+\delta^2}.
\label{eq:Wpm}
\end{equation}
\(c_k\) is the per-packet coupling and \(\Gamma_k\) the per-packet homogeneous width
(\cref{sec:profiles}); by default \(c_k\) is fixed by the peak selectivity and
\(\Gamma_k=\Gamma\). The RF currents are
\begin{equation}
J^+_k=W^+_k(n_{k+}-n_{k0}),\quad J^-_k=W^-_k(n_{k0}-n_{k-}),
\label{eq:rfcurrents}
\end{equation}
\begin{equation}
\dot n_{k+}\big|_{\rm RF}=-J^+_k,\qquad
\dot n_{k0}\big|_{\rm RF}=J^+_k-J^-_k,\qquad
\dot n_{k-}\big|_{\rm RF}=J^-_k .
\label{eq:rfpops}
\end{equation}
There is no row normalization, no edge correction and no \(\Delta R\) factor: \(U_j\) is a
carrier power, not a density. Refining the grid therefore requires the \emph{total} power of a
distributed command to be preserved. In the limit \(\Gamma\to0\), with one carrier exactly on bin
\(b\), \cref{eq:Wpm} reduces to the ideal-bin operator of the first release:
\(W^+_k\propto\delta_{kb}\) and \(W^-_k\propto\delta_{k\bar b}\).

\paragraph{Implementation.} With uniform \(\Gamma\), \(K(x_k-x_j)\) is a symmetric Toeplitz matrix.
The product \(KU\) is evaluated by FFT on a zero-padded circulant embedding of length
\(L\ge2N-1\), which avoids wrap-around. On the symmetric grid,
\(\sum_jK(-x_k-x_j)U_j=(KU)_{\kb}\), so \(W^-\) is the reversed \(W^+\) field and costs nothing
extra. The transpose needed by the optimizer's adjoint is the same convolution. The last command's
response is cached, because one command is held for many steps. With per-packet widths, the
kernels \(K^\pm\) are built as dense \(N\times N\) matrices (\code{RFRateOperator}).

\subsection{Per-spin coupling, selectivity, and mechanism profiles}
\label{sec:profiles}

\Cref{eq:Wgolden} contains no density of spins. Whether horn and pedestal spins have different
per-spin RF rates (different \(T_2\) or coupling) is a material question that present burn data
do not settle. The model therefore provides two positive per-packet factors. Both multiply the
\emph{two} transitions of a packet, so each packet keeps its mirror relationship.

\paragraph{Peak selectivity and pedestal gain.}
\begin{equation}
c_k = \Bigl[1+(1-\sigma_k)\bigl((\rho_k/\rho_{\max})^{s}-1\bigr)\Bigr]
\bigl[1+(G_{\rm p}-1)\sigma_k\bigr]\,f^{\rm rf}_k .
\label{eq:ck}
\end{equation}
\begin{itemize}
\item \(\sigma_k\) is a cubic smoothstep in \(\bigl||x_k|-R_c\bigr|\): 0 inside the peak bands
(half-width \code{rf\_peak\_half\_width\_R} around \(\pm R_c\)) and 1 beyond the taper
\code{rf\_transition\_width\_R}.
\item \(s\) is \code{rf\_peak\_selectivity}. The densest packet and the outer pedestal keep unit
coupling.
\item \(G_{\rm p}\) is \code{rf\_coupling\_pedestal\_gain}.
\end{itemize}

\paragraph{General mechanism profiles.}
Every mechanism \(X\in\{\)RF coupling, RF width, exchange, DNP, \(T_1\}\) can carry the factor
\begin{equation}
f^X_k = D^X_k\sum_{s\in\{\rm CD,OD\}} w_s(x_k)\,m^X_s\,\bigl[1+a^X_2P_2(\cos\theta_s)+a^X_4P_4(\cos\theta_s)\bigr],
\label{eq:profile}
\end{equation}
where:
\begin{itemize}
\item \(m^X_{\rm CD}=1\) and \(m^X_{\rm OD}\) is \code{X\_od\_factor};
\item \(D^X_k=\clip\bigl((\rho_k/\bar\rho_\mu)^{e_X},c^{-1},c\bigr)\), renormalized to
\(\mu\)-weighted mean one;
\item \(P_2\) and \(P_4\) are Legendre polynomials.
\end{itemize}
The angular form is the natural one for a coupling or a \(T_2\) that depends on orientation, for
example through the dipolar local field or a motional spectral density at the EFG angle. All
defaults give \(f^X\equiv1\). The RF width enters as \(\Gamma_k=\Gamma f^{\rm width}_k\) on the
\emph{receiving} packet. At fixed \(\kappa\), a narrower packet has a higher on-resonance rate
\(\kappa/\Gamma_k\) and a narrower hole. This is the physically motivated route to
horn-favouring selectivity. The \code{d\_butanol\_basic} and \code{d\_butanol\_extend}
materials use it (\cref{sec:materials}).

These profiles implement the separately calibrated RF, recovery and DNP weights
\(w^{\rm RF}_k\), \(w^{\rm rec}_k\), \(w^{\rm DNP}_k\) that the first-release note proposed as a
future extension.

\paragraph{Legacy RF modes.} \code{rf\_response\_mode = legacy} reproduces the first-release law,
in which RF is multiplied by the DNP density weight. \code{independent} adds a separate density
exponent and a smooth pedestal gain; it is the MaterialTuning-V1 law. Both use one-bin rates,
\(W^+_k=w_kU_k\) and \(W^-_k=w_kU_{\kb}\). They are kept to reproduce earlier results; the
shipped materials use \code{spectral}.

\subsection{How RF changes \texorpdfstring{\(P\)}{P} and \texorpdfstring{\(Q\)}{Q}}
\label{sec:rf-PQ}

From \cref{eq:Pdef,eq:Qdef,eq:rfpops},
\begin{equation}
\dot P\big|_{\rm RF} = -\sum_k\bigl(J^+_k+J^-_k\bigr),\qquad
\dot Q\big|_{\rm RF} = 3\sum_k\bigl(J^-_k-J^+_k\bigr)=3\sum_k\bigl(W^-_k\Delta^-_k-W^+_k\Delta^+_k\bigr),
\label{eq:dQrf}
\end{equation}
with \(\Delta^+_k=n_{k+}-n_{k0}\) and \(\Delta^-_k=n_{k0}-n_{k-}\). For the usual monotone
spin-temperature populations, RF reduces \(|P|\). It raises \(Q\) exactly when it saturates more
\(\Iminus\) population difference than \(\Iplus\) population difference. For one narrow carrier
at \(R\), the sign of \(\dot Q\) is the sign of
\(c_{\overline{k_+}}\Delta^-_{\overline{k_+}}-c_{k_+}\Delta^+_{k_+}\). In words, the
coupling-weighted \(\Iminus\) signal at that frequency must exceed the \(\Iplus\) signal. At
\(P>0\) this selects two regions:
\begin{itemize}
\item the little peak (\(R\approx-0.93\), the \(\Iminus\) horn);
\item the region \(R\gtrsim1.1\), where only the \(\Iminus\) line exists.
\end{itemize}
The optimizer's candidate mask (\cref{sec:candidates}) is this criterion evaluated with the full
kernel.

\subsection{Direct holes, mirror peaks, and the one-half gain rule}
\label{sec:mirror}

Using \cref{eq:displayIplus,eq:displayIminus}, take the ideal-bin limit: a carrier on bin \(b\)
with \(\Gamma\to0\), so that only \(J^+_{k_+}\) acts on packet \(k_+\) and only \(J^-_{k_-}\) on
packet \(k_-\). Write \(\mathcal K=\mathcal C_{\rm disp}/\Delta R\). Then
\begin{align}
  \frac{\dd\Iplus(R_b)}{\dd t}\bigg|_{\rm RF}
  &=-2\mathcal K J^+_{k_+},
  &\frac{\dd\Iminus(-R_b)}{\dd t}\bigg|_{\rm RF}
  &=+\mathcal K J^+_{k_+},
  \label{eq:directplus}\\
  \frac{\dd\Iminus(R_b)}{\dd t}\bigg|_{\rm RF}
  &=-2\mathcal K J^-_{k_-},
  &\frac{\dd\Iplus(-R_b)}{\dd t}\bigg|_{\rm RF}
  &=+\mathcal K J^-_{k_-},
  \label{eq:directminus}
\end{align}
and therefore
\begin{equation}
  \dot\Iminus(-R_b)=-\frac12\dot\Iplus(R_b),\qquad
  \dot\Iplus(-R_b)=-\frac12\dot\Iminus(R_b).
  \label{eq:halfgain}
\end{equation}
The integrated mirror gain is one half of the direct loss \cite{ClementKeller2023}. The mirror
peaks do not contradict depolarization: each direct transition loses twice the signal gained in
the mirror projection, leaving a net reduction in branch area,
\(\dot P|_{\rm RF}=-J^+_{k_+}-J^-_{k_-}\).

\paragraph{Finite linewidth.} With \(\Gamma>0\), the two conditions behind the one-half rule
break:
\begin{itemize}
\item a carrier near the line centre (\(|R|\lesssim\Gamma\)) drives \emph{both} transitions of the
same packet;
\item the carrier's Lorentzian wings reach neighbouring packets.
\end{itemize}
The ratio is then no longer fixed, which a regression test checks. Spin diffusion,
DNP and \(T_1\) also act during and after a burn. A measured mirror/direct ratio that differs from
\(\tfrac12\) is therefore a diagnostic of these processes, not a contradiction of the bookkeeping.

%==================================================================================================
\section{Spectral spin diffusion}
\label{sec:diffusion}

\subsection{Flip-flops between quadrupole-split spin-1 nuclei}
\label{sec:flipflop}

The secular dipolar coupling between two like spins contains the flip-flop operator
\(\hat I_+^{(1)}\hat I_-^{(2)}+\mathrm{h.c.}\). For spin 1 it connects three kinds of pair
states \cite{Abragam1961,SuterErnst1985,Eckman1982}. Using \cref{eq:levels-rot}, the energy
mismatch each process must exchange with the dipolar bath, when packet \(i\) changes first, is:
\begin{align}
\text{same-branch }(+):&\quad |{+}_i0_j\rangle\leftrightarrow|0_i{+}_j\rangle, & \delta&=x_i-x_j,\\
\text{same-branch }(-):&\quad |0_i{-}_j\rangle\leftrightarrow|{-}_i0_j\rangle, & \delta&=-(x_i-x_j),\\
\text{cross-branch}:&\quad |0_i0_j\rangle\leftrightarrow|{+}_i{-}_j\rangle, & \delta&=x_i+x_j .
\end{align}
The squared matrix elements are equal:
\begin{equation}
|\langle0{+}|\hat I_+^{(2)}\hat I_-^{(1)}|{+}0\rangle|^2=|\langle{+}{-}|\hat I_+^{(1)}\hat I_-^{(2)}|00\rangle|^2=4 .
\end{equation}
For a spatially random dipolar network, the cross-branch to same-branch ratio
\(r_{\rm x}\) (\code{cross\_branch\_ratio}) is therefore 1, which is the default. Same-branch
exchange needs \(x_i\approx x_j\); cross-branch exchange needs \(x_i\approx-x_j\), so that the
\(\Iplus\) line of one packet overlaps the \(\Iminus\) line of the other.

A flip-flop with mismatch \(\delta\) proceeds at a rate proportional to the zero-quantum (ZQ)
spectral density at \(\delta\). That density is modelled as a normalized-to-one overlap,
\begin{equation}
S(\delta)=\frac{1}{1+(\delta/w)^2}\quad\text{(Lorentzian, default)}\qquad\text{or}\qquad
S(\delta)=e^{-\delta^2/2w^2},
\label{eq:overlap}
\end{equation}
with width \(w\) given by \code{zq\_width\_R}. A Lorentzian follows from an exponentially
decaying ZQ correlation function. Its algebraic tails allow weak exchange far off resonance,
with the energy taken by the dipolar bath. Outside the extended engine, that bath is treated as
infinitely hot and infinitely large. \(w\) is \emph{not} the static line width \(\gamma\), nor
the RF homogeneous width \(\Gamma\). A physical estimate is \(w\approx2\Gamma\), the convolution
of two homogeneous lines.

\subsection{Pair currents}
\label{sec:paircurrents}

Treat each packet as well mixed, and assume that the spatial neighbourhood of a spin does not
depend on its EFG orientation (\code{orientation\_corr\_fraction} \(=0\)). The forward rate of
\((+_i,0_j)\to(0_i,+_j)\) is then \(K_0\mu_i\mu_jS_{ij}\,p_{i+}p_{j0}\), where
\(S_{ij}=S(x_i-x_j)\) and \(K_0\) is \code{diffusion\_scale}. The net currents are of mass-action
form:
\begin{align}
J^{(+)}_{ij}&=K_0 S_{ij}\bigl(n_{i0}n_{j+}-n_{i+}n_{j0}\bigr),\label{eq:Jsp}\\
J^{(-)}_{ij}&=K_0 S_{ij}\bigl(n_{i-}n_{j0}-n_{i0}n_{j-}\bigr),\label{eq:Jsm}\\
J^{({\rm x})}_{ij}&=r_{\rm x}K_0 X_{ij}\bigl(n_{i0}n_{j0}-n_{i+}n_{j-}\bigr),\qquad
X_{ij}=S(x_i+x_j)\bigl(1-\tfrac12\delta_{ij}\bigr).\label{eq:Jx}
\end{align}
The factor \(\tfrac12\) on the diagonal of \(X\) counts unordered pairs drawn from one packet; it
matters only near \(R=0\). The packet updates are
\begin{align}
\dot n_{i+}\big|_{\rm SQ+X} &= \textstyle\sum_j J^{(+)}_{ij}+\sum_jJ^{({\rm x})}_{ij},\notag\\
\dot n_{i-}\big|_{\rm SQ+X} &= \textstyle-\sum_j J^{(-)}_{ij}+\sum_jJ^{({\rm x})}_{ji},\label{eq:diffpops}\\
\dot n_{i0}\big|_{\rm SQ+X} &= -\dot n_{i+}-\dot n_{i-}.\notag
\end{align}
The capacities appear exactly once, through \(n=\mu p\). The Pake density therefore sets the
number of \emph{partners} a packet has within the ZQ width, and is never applied again as a rate
multiplier.

\paragraph{Microwave factor and per-packet exchange profiles.} With DNP on, \(K_0\) is
multiplied by \code{microwave\_diffusion\_factor} (default 1), which allows for diffusion
modified near paramagnetic centres. The exchange profile of \cref{eq:profile} enters as
\(S_{ij}\to\sqrt{f^{\rm ex}_if^{\rm ex}_j}\,S_{ij}\). This keeps every kernel symmetric and every
current reversible.

\paragraph{Orientation correlation.} If the real-space neighbourhood correlates with EFG
orientation, the kernel is multiplied by
\((1-f_\theta)+f_\theta\,e^{-(\theta_i-\theta_j)^2/2\sigma_\theta^2}\). This is optional and off
in all shipped materials.

\subsection{Double-quantum transport}
\label{sec:dq}

The process \((+_i,-_j)\leftrightarrow(-_i,+_j)\) exchanges vector polarization between packets.
It changes neither \(N^\pm\) nor \(E_{\rm spin}\) (\cref{eq:levels-rot}), so it needs no
single-quantum frequency matching. Microscopically it is a second-order process: two flip-flops
through a virtual intermediate state, whose amplitude involves the mean squared coupling. It is
folded into one effective ratio \(r_{\rm DQ}\) (\code{double\_quantum\_ratio}) to \(K_0\):
\begin{equation}
J^{\rm DQ}_{ij}=r_{\rm DQ}K_0\,\sqrt{f^{\rm ex}_if^{\rm ex}_j}\,C_{ij}\bigl(n_{i-}n_{j+}-n_{i+}n_{j-}\bigr),\qquad
\dot n_{i+}\mathrel{+}=\textstyle\sum_jJ^{\rm DQ}_{ij},\quad
\dot n_{i-}\mathrel{-}=\textstyle\sum_jJ^{\rm DQ}_{ij},
\label{eq:Jdq}
\end{equation}
where \(C_{ij}\) is a spatial/orientation correlation matrix, unity off the diagonal by default.
For \(C\equiv1\) the sum separates into \(\sum_jn_{j\pm}\), so the cost is \(O(N)\). In the
fitted d-butanol material (\(r_{\rm DQ}=3.03\), broad \(w=0.717\)), DQ is a dominant refill
channel. Refilling is then largely nonlocal in frequency (\cref{sec:materials}).

\subsection{Conservation laws, H-theorem, and the stationary state}
\label{sec:stationary}

\paragraph{Invariants.}
\begin{itemize}
\item Each same-branch event conserves \(N^+\), \(N^0\) and \(N^-\), and hence \(P\) and \(Q\).
\item A cross-branch event \((0,0)\to(+,-)\) changes \((N^+,N^0,N^-)\) by \((+1,-2,+1)\). It
conserves \(P\) and changes \(Q\) by \(+6\) per event, in units of the packet population.
\item DQ events conserve \(N^\pm\).
\item All channels conserve every \(\mu_k\).
\end{itemize}
So the internal dynamics never changes \(P\), and it changes \(Q\) only through the cross-branch
channel.

\paragraph{H-theorem.} Write each channel as \(\sum_{\rm pairs}\kappa\,(a-b)\) with
\(\kappa\ge0\), where \(a\) and \(b\) are the forward and backward products of fractions. Then
\begin{equation}
\frac{\dd}{\dd t}\sum_{k,m}n_{km}\ln p_{km} = \sum_{\rm pairs}\kappa\,(a-b)\ln\frac{a}{b}\;\ge\;0,
\label{eq:Htheorem}
\end{equation}
so a relative entropy is a Lyapunov function, and stationary states satisfy detailed balance
pair by pair:
\begin{itemize}
\item same-branch: \(p_{i+}/p_{i0}\) and \(p_{i-}/p_{i0}\) are equal for all packets connected by
\(S>0\);
\item cross-branch: \(p_0^2=p_+p_-\);
\item DQ: \(p_{i+}/p_{i-}\) is uniform.
\end{itemize}
On a connected graph with \(r_{\rm x}>0\), these conditions force \(p_{km}\propto r^m\) with one
\(r\) for all packets, which is a Boltzmann state. Since \(P\) is conserved, the stationary state
is unique: \emph{the Boltzmann line at the remaining \(P\)}.

The first release imposed this behaviour through a reference state \(\bm n^{\rm B}[P(t)]\) and a
projection. Here it is a theorem about the currents. The tests check it directly: stationarity of
every channel at Boltzmann, positive entropy production, relaxation of arbitrary states to the
unique Boltzmann line at the remaining \(P\), and non-stationarity of uniform non-Boltzmann
states.

\paragraph{Restricted channels.} Without the cross channel (\(r_{\rm x}=0\)), every state with
uniform fractions is stationary, including non-Boltzmann ones, and \(Q\) is frozen at fixed \(P\).
This is the \code{d\_butanol\_same\_branch\_transport} hypothesis. It is physically extreme, but
it is the regime in which steady RF\,+\,DNP gives the largest tensor enhancement
(\cref{sec:regimes}).

\subsection{Fast evaluation: Toeplitz and Hankel products}
\label{sec:fft}

On the symmetric uniform grid, \(S(x_i-x_j)\) depends only on \(i-j\), so it is Toeplitz. Because
\(-x_j=x_{\bar\jmath}\),
\begin{equation}
\sum_j S(x_i+x_j)\,v_j=\sum_j S(x_i-x_{\bar\jmath})\,v_j=\bigl(S\,v^{\rm rev}\bigr)_i,
\end{equation}
so the cross kernel is the same Toeplitz product applied to the reversed vector, minus half the
diagonal. Expanding the mass-action currents gives, for example,
\begin{equation}
\textstyle\sum_jJ^{(+)}_{ij}=K_0\bigl[n_{i0}(Sn_+)_i-n_{i+}(Sn_0)_i\bigr],
\end{equation}
and the \(i=j\) terms cancel identically. All three population columns are convolved in one
batched real FFT. \code{Spin1Model.\_spin\_diffusion\_terms} provides three algebraically
identical evaluators, and the tests require agreement to rounding:
\begin{itemize}
\item FFT, the default with no orientation correlation and no ambiguous hard cutoff;
\item sparse CSR matrices;
\item explicit pair lists with \code{np.add.at}.
\end{itemize}
The default keeps the full Lorentzian tails on the finite grid; a positive
\code{kernel\_cutoff\_widths} gives a banded kernel.

%==================================================================================================
\section{DNP and spin-lattice relaxation}
\label{sec:dnp}

\subsection{DNP source term}

DNP drives each packet toward the Boltzmann fractions at a chosen saturation polarization:
\begin{equation}
  \dot{\bm n}_k\big|_{\rm DNP}
  =k_{\rm D}\,w^{\rm D}_k\left(\mu_k\,\bm p^{\rm B}(P_{\rm sat})-\bm n_k\right),
  \qquad
  w^{\rm D}_k=\frac{\clip\bigl((\rho_k/\bar\rho_\mu)^{\alpha_{\rm D}},c^{-1},c\bigr)}
  {\sum_j\mu_j\clip\bigl((\rho_j/\bar\rho_\mu)^{\alpha_{\rm D}},c^{-1},c\bigr)}\,f^{\rm dnp}_k,
  \label{eq:dnpterm}
\end{equation}
with \(\alpha_{\rm D}=\) \code{capacity\_rate\_power} and \(c=\) \code{capacity\_rate\_clip}. Because
the target and the packet each sum to \(\mu_k\), DNP redistributes populations within a packet
without changing its spin count. With RF and \(T_1\) absent, every packet approaches the same
target fractions, and the total polarization approaches \(P_{\rm sat}\) without overshoot. The
density weight makes horn packets approach the target faster during the transient. The global
build-up is then multi-exponential, with
\begin{equation}
\dot P\big|_{\rm DNP}=k_{\rm D}\sum_kw^{\rm D}_k\bigl\{\mu_k[p^{\rm B}_+-p^{\rm B}_-](P_{\rm sat})-(n_{k+}-n_{k-})\bigr\}.
\end{equation}

\paragraph{Physical reading.} In the thermal-mixing regime, the nuclear Zeeman system reaches a
common spin temperature with the electron non-Zeeman reservoir
\cite{Borghini1968,AbragamGoldman1978,Wenckebach2016}. The nuclear populations then approach a
Boltzmann distribution that is the same at every orientation of a powder. \Cref{eq:dnpterm} is
the simplest closure with that fixed point and one build-up rate. It does not separate the solid
effect, the cross effect, thermal mixing or the differential solid effect, and it does not
represent nuclear-spin diffusion barriers around paramagnetic centres. For d-butanol the fit to
the 2014 build-up gives \(k_{\rm D}=2.30\times10^{-3}\)\,s\(^{-1}\) (\(\tau\approx435\)\,s) and
\(P_{\rm sat}=0.476\) (held-out RMS 0.79\%).

\subsection{Optional longitudinal relaxation}

\begin{equation}
  \dot{\bm n}_k\big|_{T_1}
  =k_1f^{T_1}_k\left[\mu_k\bm p^{\rm B}(P_{\rm eq})-\bm n_k\right],
  \label{eq:T1term}
\end{equation}
with an optional per-packet profile \(f^{T_1}_k\), which allows an orientation-dependent
\(T_1(\theta)\) \cite{Keller2026}. The term is off in the shipped materials (\(k_1=0\)), so
no-DNP simulations isolate RF depolarization and internal redistribution.

%==================================================================================================
\section{Extended engine: site-resolved populations and the dipolar reservoir}
\label{sec:extended}

\code{ExtendedCore} is used when \code{site\_resolved} or \code{dipolar\_enabled} is set. It
requires spectral RF, no orientation correlation and no kernel cutoff. With one site and
\(\beta=0\) it reproduces the basic equations and their adjoint exactly, which is tested.

\subsection{Site resolution}

The state becomes \(n_{s,k,m}\), with \(s\in\{\mathrm{CD},\mathrm{OD}\}\), and capacities
\begin{equation}
\mu_{s,k}=\mu_k\,\frac{a_s(x_k)}{\sum_{s'}a_{s'}(x_k)},
\end{equation}
where \(a_s\) is the site absorption in \cref{eq:multisite}. Each site has its own orientation
map, so it has its own profiles \(f^X_{s,k}\) (\cref{eq:profile} without the site sum) and its
own RF operator. Exchange between sites is scaled by a \(2\times2\) mixing matrix \(M\) with
off-diagonal elements \code{exchange\_intersite\_factor}:
\begin{equation}
(\mathcal{H}V)_{s} = \sqrt{f_s}\sum_t M_{st}\,S*\bigl(\sqrt{f_t}\,V_t\bigr).
\end{equation}
The public state is the site sum \(n_k=\sum_sn_{s,k}\). If \(n\) is edited externally, it is
redistributed over sites with equal level fractions. AFP acts on each site exactly.

\subsection{Dipolar reservoir in the high-temperature approximation}

The secular dipolar Hamiltonian \(\hat H_D\) is treated as a separate thermodynamic reservoir with
inverse temperature \(\beta\), in \(1/R\) units \cite{Provotorov1962,Goldman1970,AbragamGoldman1982}.
To first order in \(\beta\),
\begin{equation}
E_d=\langle\hat H_D\rangle=-C_d\,\beta,\qquad C_d = c_d\sum_{s,k}\mu_{s,k}x_k^2,
\end{equation}
where \(c_d\) is \code{dipolar\_capacity\_ratio}; it relates the dipolar heat capacity to the
second moment of the offset distribution. Consider a process that takes energy
\(\epsilon_{\rm in}\) from the reservoir. Detailed balance requires
\(W_{\rm f}/W_{\rm b}=e^{-\beta\epsilon_{\rm in}}\). Expanding to first order, with
\(W_{\rm f,b}=W(1\mp\tfrac12\beta\epsilon_{\rm in})\), gives the linearized current
\begin{equation}
J = W\Bigl[(n_a-n_b)-\beta\,\epsilon_{\rm in}\,\tfrac12(n_a+n_b)\Bigr].
\label{eq:dbalance}
\end{equation}
Only the reservoir factor is linearized; the Zeeman populations are not.

\paragraph{RF.} A carrier at \(x_j\) drives the \(\Iplus\) transition of packet \(k\). The step
\(+\to0\) raises \(E_{\rm spin}\) by \(x_k\) and the photon supplies \(x_j\), so the reservoir
supplies \(x_k-x_j\). Summing over carriers turns the expansion into moments of the kernel:
\begin{align}
J^+_k &= G^+_k\Delta^+_k - \beta H^+_k\,\tfrac12(n_{k+}+n_{k0}),\\
G^+_k &= c_k\textstyle\sum_jK_{kj}U_j,\qquad
H^+_k=c_k\textstyle\sum_jK_{kj}(x_k-x_j)U_j,\qquad
\mathcal Q^+_k=c_k\textstyle\sum_jK_{kj}(x_k-x_j)^2U_j .
\end{align}
The energy given to the reservoir is
\(-\sum_k\bigl(H^+_k\Delta^+_k-\beta\mathcal Q^+_k\tfrac12(n_{k+}+n_{k0})\bigr)\). The \(\Iminus\)
terms are analogous, with offset \(-x_k-x_j\). The moments are evaluated as
\(x\circ(KU)-K(xU)\) and \(x^2\circ(KU)-2x\circ K(xU)+K(x^2U)\): three convolutions. Off-resonant
irradiation cools or heats the reservoir with the Provotorov sign, which is tested.

\paragraph{Flip-flops.} Same-branch exchange carries the mismatch \(x_i-x_j\):
\begin{equation}
J^{(+)}_{ij}=K_0S_{ij}\Bigl[(n_{i0}n_{j+}-n_{i+}n_{j0})+\tfrac12\beta\xi(x_i-x_j)(n_{i0}n_{j+}+n_{i+}n_{j0})\Bigr],
\end{equation}
which is evaluated with \(x\circ(Sn)-S(xn)\) products. Cross-branch exchange carries
\(x_i+x_j\). DQ exchange is energy-neutral and gets no \(\beta\) term. \(\xi\)
(\code{dipolar\_exchange\_coupling}) scales the flip-flop part; \(\xi=0\) leaves only the
RF--Provotorov coupling.

\paragraph{Reservoir equation.}
\begin{equation}
\dot\beta = -\frac{\dot E_{\rm released}}{C_d}-r_{1d}\,(\beta-\beta_L),
\end{equation}
where \(\dot E_{\rm released}=-\dot E_{\rm spin}\) from exchange, plus the RF mismatch terms. With
exchange alone, \(E_{\rm spin}+E_d\) is conserved \emph{exactly} by the right-hand side, not just
to first order; the tests check this to rounding. In the \code{d\_butanol\_full\_physics}
material, \(c_d=0.01\), and \(|\beta\delta|\) reaches 0.2--0.7 in slow, power-limited runs.
Quantitative use in that regime needs the exponential form of \cref{eq:dbalance}. For uncapped
single-burn optima, site resolution and the reservoir changed \(Q\) by \(\le0.0003\).

%==================================================================================================
\section{Adiabatic fast passage}
\label{sec:afp}

An adiabatic frequency sweep across a window inverts every transition resonant in it
\cite{GarwoodDelaBarre2001}. The window at \(x_k\) contains the \(\Iplus\) transition of packet
\(k\) and the \(\Iminus\) transition of packet \(\kb\). The sweep is modelled as an instantaneous
swap with efficiency \(\varepsilon\):
\begin{align}
(n_{k+},n_{k0})&\to\bigl(\varepsilon n_{k0}+(1-\varepsilon)n_{k+},\;\varepsilon n_{k+}+(1-\varepsilon)n_{k0}\bigr),\\
(n_{\kb0},n_{\kb-})&\to\bigl(\varepsilon n_{\kb-}+(1-\varepsilon)n_{\kb0},\;\varepsilon n_{\kb0}+(1-\varepsilon)n_{\kb-}\bigr).
\end{align}
\code{mirror\_safe\_indices} keeps at most one index from each mirror pair, so no packet is
swapped twice. \code{qneg\_sweep\_indices} selects the bins where the signed tensor signal
\(\Iplus-\Iminus\) is negative (the GUI's ``optimal profile'' region). The swap conserves
\(\mu_k\). Its finite duration, and any relaxation during it, are neglected. In the
single-burn studies, a two-window AFP reached \(Q\approx0.298\) from \(P_0=0.50\) at full
efficiency.

%==================================================================================================
\section{Physical regimes represented by the model}
\label{sec:regimes}

\subsection{DNP off, RF on}

With \(k_{\rm D}=0\) and \(k_1=0\), RF is the only term that changes vector polarization. Exchange
feeds population difference into the addressed packets from the rest of the line, while RF
equalizes them. The depolarization rate is therefore controlled by four things:
\begin{itemize}
\item the RF power and linewidth;
\item the two per-spin couplings;
\item the two current transition differences;
\item the rate at which exchange replenishes those differences.
\end{itemize}
The replenishment rate depends on the number of partners within the ZQ width (\(\propto(S\mu)_k\))
and on the DQ partner sum (\cref{sec:peak-pedestal}).

\subsection{DNP off, RF switched off after a burn}

Once RF is off, \(\dot P=0\). The exchange channels remove the hole and mirror structure while
conserving the reduced \(P\). The line relaxes to the Boltzmann line at \(P_{\rm after\,burn}\),
not the original \(P_0\) (\cref{sec:stationary}). Spectral recovery is not repolarization. The
tensor polarization relaxes toward \(Q_{\rm B}(P_{\rm after\,burn})\) only as fast as
cross-branch exchange allows.

\subsection{DNP on, RF off}

The line builds toward \(P_{\rm sat}\). Packets have different transient time constants because
of \(w^{\rm D}_k\), but the common target fractions give the same asymptotic Boltzmann shape. The
transient can be non-Boltzmann across \(R\), because the horns reach the target faster than
weak regions. Exchange smooths that difference.

\subsection{DNP on, RF on}

The system approaches a driven steady state, \(\dot{\bm n}=0\) in \cref{eq:totalrate}. DNP
rebuilds population difference, exchange transports it through packet space, and RF removes it.
Hole depth and \(Q\) are set by that balance, not by RF power alone. This regime depends most on
the model:
\begin{itemize}
\item the best steady \(Q\) under continuous DNP is 0.20--0.21 with the fitted mixed transport;
\item it is 0.24 with same-branch transport only, and 0.30 at hundredfold same-branch exchange;
\item the best steady burn sits where only the \(\Iminus\) line exists (\(R\approx+1.5\) to
\(+2.0\)), about 15--30\,mW;
\item burning the horn at the fitted DNP rate is harmful;
\item pulsing or bunching the RF in time never beats a steady burn, because incoherent saturation
is concave in dose.
\end{itemize}

\subsection{Pulsed operation: optimize, evolve, re-optimize}

In practice, ssRF is not run continuously. A pulse is applied, the NMR and control loop take a few
seconds to tens of seconds, and the state is measured and re-optimized. After an optimized pulse
the optimum changes:
\begin{itemize}
\item power moves from the \(\Iminus\)-only region to the little peak;
\item the pulse becomes shorter and stronger;
\item re-applying a fixed program loses up to half of the available gain.
\end{itemize}
This regime motivates the closed-loop data generator (\cref{sec:datagen}).

\subsection{Negative polarization}

All equations are signed. Reversing the spin-temperature polarization reverses the branch signs
and the RF currents. For monotone orderings RF still reduces \(|P|\). DNP can build toward a
negative \(P_{\rm sat}\), and \(T_1\) can relax toward either sign. General non-Boltzmann states
require evaluating the currents explicitly, and the optimizer and candidate rule use the signed
populations directly.

%==================================================================================================
\section{Numerical implementation}
\label{sec:numerics}

\subsection{Grid and state initialization}

\begin{center}
\begin{tabular}{lll}
\toprule
Material & Grid & Step \(\delta t\) \\
\midrule
d-butanol materials & \(N=500\), \(R\in[-6.186,6.186]\), \(\Delta R=0.0248\) & \(0.03\)\,s \\
\code{demo\_single\_pake} & \(N=500\), \(R\in[-3,3]\), \(\Delta R=0.0120\) & \(0.03\) (arbitrary units) \\
\bottomrule
\end{tabular}
\end{center}
The initial state is \cref{eq:initpacket}. A copy is stored as the reference used for plotting
direct and mirror changes. It does not enter the dynamics.

\subsection{Time integration}
\label{sec:integration}

The live model integrates \cref{eq:totalrate} with positivity-preserving forward Euler and
adaptive substeps:
\begin{equation}
\bm n\leftarrow \bm n+h\,\dot{\bm n}(\bm n),\qquad
h=\min\Bigl(t_{\rm event}-t,\;\frac{0.2}{\lambda_{\max}},\;h_{\rm prog}\Bigr).
\label{eq:euler}
\end{equation}
\begin{itemize}
\item \(t_{\rm event}\) is the next pulse boundary, so every pulse starts and stops exactly.
\item \(h_{\rm prog}\) is an optional per-program ceiling, applied only while RF is on. The
designer sets it to the validated playback step.
\item \(\lambda_{\max}\) is a conservative bound on the outgoing rate of any population:
\(W^++W^-\), the DNP and \(T_1\) rates, and the exchange bound
\(2K_0(S\mu)_k+2r_{\rm x}K_0|X\mu|_k+r_{\rm DQ}K_0(C\mu)_k\), evaluated at full occupancy.
\end{itemize}
For the linear part, \(h\lambda\le1\) keeps every population nonnegative; the factor 0.2 leaves
margin for the bilinear terms. Within each material step \(\delta t\), the step is split as
needed. After each substep the code:
\begin{enumerate}
\item raises an error if any population is below \(-10^{-13}\mu_k\), instead of silently
clipping a real error;
\item floors at \(10^{-30}\);
\item rescales each packet to its exact mass \(\mu_k\), which removes rounding drift;
\item accumulates the delivered exposure \(\sum h\,U_j\) per bin.
\end{enumerate}

\paragraph{Why explicit Euler.} The right-hand side costs \(O(N\log N)\): three to six FFTs per
evaluation. The positivity bound sets the step anyway. And the optimizer needs a forward map whose
discrete adjoint is simple and exact (\cref{sec:adjoint}). The first-order error is controlled
\emph{a posteriori} by step halving (\cref{sec:validation}). For the fitted d-butanol material the rates are slow: with RF off, the bound allows steps of
about 3\,s, so the material step \(\delta t=0.03\)\,s sets the resolution. For the 420 optimized
programs of the chain dataset (\cref{sec:datagen}), at the 0.3\,W cap, the validated playback
step was 0.015--0.03\,s. The rate bound limits the step only for uncapped, high-power commands.

\subsection{Pulse programs and event-aware stepping}

A program is a list of \code{BinPulse}\((j,U,t_0,\tau)\), each acting on the half-open interval
\([t_0,t_0+\tau)\); profiles add where they overlap. \code{CompiledProgram} sorts all edges into
one event list and rebuilds the command vector \(U(t)\) from the pulses active in each interval,
instead of adding and subtracting increments. This prevents residual RF after the last pulse.
\code{start\_program} anchors program time to the current simulation time, and the integrator
splits at every absolute event time. If two boundaries fall below floating-point resolution at a
late simulation time, the start is refused, so a pulse is never silently merged.

\subsection{Spectral projection and display}

The \(\Iplus\) array is plotted at \(R=x_k\). The \(\Iminus\) array is plotted at \(R=-x_k\),
which is the reversed array on the symmetric grid. The total is their binwise sum. Display noise,
when enabled, is added to the plotted arrays at redraw time only; it never feeds back into the
populations.

\subsection{Conservation summary}

\begin{itemize}
  \item Every packet preserves \(\sum_mn_{km}=\mu_k\) after each step.
  \item Same-branch and DQ exchange conserve \(N^+,N^0,N^-\), and hence \(P\) and \(Q\),
  analytically.
  \item Cross-branch exchange conserves \(P\) and changes \(Q\) toward \(Q_{\rm B}(P)\).
  \item RF conserves packet spin counts and changes \(P\) and \(Q\) according to
  \cref{eq:dQrf}.
  \item DNP and \(T_1\) conserve packet spin counts and drive \(P\) and \(Q\) toward their
  targets.
  \item With DNP, \(T_1\) and RF off, \(P\) is conserved and the state relaxes to the unique
  Boltzmann line at that \(P\).
  \item In the extended engine, exchange conserves \(E_{\rm spin}+E_d\).
\end{itemize}

\subsection{Cost}

One evaluation of the right-hand side costs:
\begin{itemize}
\item one convolution for the RF field;
\item one batched three-column FFT convolution for same-branch exchange;
\item one more for cross-branch exchange;
\item \(O(N)\) for DQ, DNP and \(T_1\).
\end{itemize}
At \(N=500\) this takes tens of microseconds in NumPy. The extended engine multiplies the cost by
the number of sites, plus three convolutions per site for the reservoir moments.

%==================================================================================================
\section{Snapshots, materials, and reproducibility}
\label{sec:snapshots}

A \emph{material} file (\code{ssrf.material.v1}) stores every kinetic parameter, the optimizer
defaults and GUI settings. Older files are upgraded by inserting parameter groups at their neutral
defaults, and loading is transactional.

A \code{PopulationSnapshot} stores:
\begin{itemize}
\item the full parameter dictionary;
\item \(n\), the reference state, \(\mu\) and the grid;
\item the simulation time;
\item for the extended engine, \(n_{s,k,m}\) and \(\beta\).
\end{itemize}
Its fingerprint is a SHA-256 hash of the kinetic parameters, \(n\), \(\mu\) and the extended
state; display-only settings are excluded. \code{make\_model} rebuilds an \code{IdealBinModel}
and refuses if any of three checks fails:
\begin{itemize}
\item the capacities recomputed from the parameters match the stored ones;
\item the populations are physical;
\item the fingerprint matches.
\end{itemize}
A snapshot is thus a self-validating initial condition, and every optimizer report records the
fingerprint of the state it was computed for.

%==================================================================================================
\section{Tensor-polarization optimizer}
\label{sec:optimizer}

\subsection{Problem}

From a frozen state \(\bm n(0)\), choose carrier powers \(U\in\mathbb{R}^N\), held constant on
\([0,T)\) (a \emph{synchronized} program), and a duration \(T\), to solve
\begin{equation}
\max_{U,T}\;Q\bigl(\bm n(T)\bigr)\quad\text{s.t.}\quad 0\le U_j\le \bar U_j,\quad
\sum_jU_j\le P_{\rm cap},\quad T\sum_jU_j\le E_{\rm cap},\quad T\in[T_{\min},T_{\max}],
\label{eq:problem}
\end{equation}
with the natural dynamics as captured (DNP, exchange and \(T_1\) on or off). Here
\(P_{\rm cap}\) is the total-power limit (\(\sum U=1\) corresponds to 0.1\,W) and \(E_{\rm cap}\)
the energy limit (\(T\sum U=34\) corresponds to 3.4\,J). The per-bin ceilings \(\bar U_j\) are
zero outside the candidate set and outside any user regions. At fixed \(T\) the two caps merge
into one linear constraint,
\begin{equation}
\sum_jU_j\le B(T)/T,\qquad B(T)=\min\bigl(E_{\rm cap},\,P_{\rm cap}T\bigr)
\label{eq:budget}
\end{equation}
(\code{effective\_exposure\_limit}). The objective is \(Q(T)\) itself. The report always
includes the no-RF endpoint, so that natural build-up is never credited to the RF.

\subsection{Discretize, then optimize}

For a given \(T\), the code fixes a step count \(K\) and step \(h=T/K\) from a bound that does
\emph{not} depend on the trial powers:
\begin{equation}
h=\min\Bigl(h_{\rm search},\,\frac{0.18}{\lambda_{\max}\bigl(\min(\bar U,B/T)\bigr)}\Bigr)
\end{equation}
(\code{prediction\_step\_plan}). The forward map
\(\Phi_T:U\mapsto \bm n_K\), with \(\bm n_{k+1}=\bm n_k+h\mathcal F(\bm n_k;U)\), is then one
fixed, smooth function of \(U\). If \(K\) depended on \(U\), the objective would be piecewise
defined, and its gradient would not exist at the switching points. \code{FrozenDynamics} holds an
algebraically regrouped copy of the right-hand side, which the tests check against
\code{Spin1Model.derivative}.

\subsection{Discrete adjoint}
\label{sec:adjoint}

Let \(\lambda_k=\partial Q/\partial\bm n_k\). Then \(\lambda_K=(1,-2,1)\) in every packet, and
\begin{equation}
\lambda_k=\lambda_{k+1}+h\Bigl(\frac{\partial\mathcal{F}}{\partial \bm n}\Big|_{\bm n_k}\Bigr)^{\!\top}\lambda_{k+1},
\qquad
\nabla_UQ = h\sum_{k=0}^{K-1}\Bigl(\frac{\partial\mathcal{F}}{\partial U}\Big|_{\bm n_k}\Bigr)^{\!\top}\lambda_{k+1}.
\label{eq:adjoint}
\end{equation}
This is the exact gradient of the discrete map \cite{GilesPierce2000,Plessix2006}. Each term
supplies a vector--Jacobian product (VJP).

\paragraph{RF.} With \(v^+=\lambda^+-\lambda^0\) and \(v^-=\lambda^0-\lambda^-\), the state
part is
\(g^+\mathrel{-}=W^+v^+\), \(g^0\mathrel{+}=W^+v^+-W^-v^-\), \(g^-\mathrel{+}=W^-v^-\), and the
command part is
\begin{equation}
\nabla_U \mathrel{+}= K^{\top}\!\bigl(c\circ(-\Delta^+)\circ v^+\bigr)
+\bigl[K^{\top}\!\bigl(c\circ(-\Delta^-)\circ v^-\bigr)\bigr]^{\rm rev}.
\end{equation}
Since \(K\) is symmetric, this is one more convolution (\code{RFRateOperator.transpose}).

\paragraph{Bilinear exchange.} Every exchange current has the form \(a_i(\mathsf{S}b)_i\) with a
symmetric kernel \(\mathsf{S}\). Its VJP contributes \(\lambda_i(\mathsf{S}b)_i\) to
\(\partial a\), and \(\mathsf{S}(\lambda\circ a)\) to \(\partial b\). The adjoint therefore uses
the same FFTs as the forward pass. The extended engine has the same structure, with the moment
kernels \(x\circ\mathsf{S}-\mathsf{S}\circ x\) and the scalar adjoint of \(\beta\) carried along.

\paragraph{Cost.} The backward pass costs about twice the forward pass. It stores the trajectory
(\(8\cdot3NK\) bytes, capped at 512\,MiB). One objective-plus-gradient evaluation therefore costs
about three forward integrations, independent of \(N\).

\subsection{Candidate bins and seeds}
\label{sec:candidates}

At \(t=0\), \cref{eq:dQrf} gives the initial slope
\begin{equation}
\frac{\partial\dot Q}{\partial U_j}\Big|_{t=0}=\Bigl[K^\top\!\bigl(-3c\circ\Delta^+\bigr)+\bigl(K^\top(3c\circ\Delta^-)\bigr)^{\rm rev}\Bigr]_j .
\end{equation}
The default candidate set (\code{union}) is the bins where this slope is positive, together with
the bins where the signed tensor signal \(\Iplus-\Iminus\) is negative. Other bins get
\(\bar U_j=0\). The seeds are:
\begin{itemize}
\item \(\bar U\) on the positive-slope bins, for the first duration;
\item the previous duration's solution, rescaled to the budget (continuation), for later
durations;
\item random fractions of \(\bar U\), for any additional starts.
\end{itemize}

\subsection{Inner solver}

\begin{itemize}
\item \textbf{Budget not binding} (\(B/T\ge\sum\bar U\)): L-BFGS-B \cite{Byrd1995} on box
constraints, in coordinates scaled to \([0,1]\).
\item \textbf{Budget binding}: SLSQP \cite{Kraft1988} with the linear inequality of
\cref{eq:budget}. The coordinate scales are capped at ten times the mean available power per bin;
otherwise per-bin ceilings far above the budget cause oversized steps. The final point is
projected onto the budget and re-evaluated, so an infeasible \(Q\) is never reported.
\end{itemize}
The objective is scaled by \(10^3\), so that small per-bin gradients are not mistaken for
convergence. The code keeps the best \emph{feasible} iterate, not the solver's last point.

\subsection{Duration scan and selection}

\begin{enumerate}
\item Solve on a geometric grid of \code{duration\_samples} values in \([T_{\min},T_{\max}]\).
\item Refine \code{duration\_refinements} times. Each refinement places \(\sqrt{T_aT_b}\) in the
wider, in \(\log T\), of the two intervals next to the best horizon. Refining both sides avoids
missing an interior optimum.
\item Re-evaluate every completed horizon at the finer endpoint step
\(\min(\delta t,h)\).
\item Select the \emph{shortest} \(T\) whose \(Q\) is within
\(\tau=\max(\tau_{\rm abs},\tau_{\rm rel}(Q_{\rm best}-Q_0))\) of the best; the defaults are
\(10^{-6}\) and 0.2\% of the gain. Shorter pulses deposit less energy and are less sensitive to
model error.
\item If no horizon beats \(Q_0\), return ``no RF''. The designer never proposes irradiation that
lowers \(Q\).
\end{enumerate}

\subsection{Validation, controls, and diagnostics}
\label{sec:validation}

\begin{itemize}
\item \textbf{Replay.} The selected program is turned into a real \code{PulseProgram} and
replayed through the unchanged scheduler and model.
\item \textbf{Half-step check.} The replay endpoint is compared with an evaluation at half the
step. For first-order Euler the full-step error is about twice that difference. The playback
step is halved, up to \code{max\_validation\_halvings} times, until \(2|\Delta Q|\le\epsilon\)
(\(2\times10^{-5}\) by default) and the state difference meets \(10\epsilon\). Otherwise the
status is \code{needs\_finer\_time\_step}. The validated step is exported as
\code{recommended\_playback\_dt}.
\item \textbf{Control.} The no-RF endpoint at the same \(T\) is reported, so the RF advantage
\(Q-Q_{\rm noRF}\) is explicit.
\item \textbf{First-order gap.} At the solution, the linearized problem
\(\max_{u\in\mathcal{U}}\nabla Q\cdot(u-U)\) over the box plus budget is a continuous knapsack.
It is solved exactly by sorting the gradient. Its value bounds the gain still available to first
order. It is a local KKT diagnostic, not a global certificate.
\item \textbf{Optional smoothing.} Gaussian smoothing of \(U\) (preserving total power and the
mask) is accepted only if, after the same replay and time-step validation, it loses no more \(Q\)
than the selection tolerance that remains unspent.
\end{itemize}

\paragraph{Scope.} This is local search over \emph{synchronized} programs from a few seeds, with
no global certificate:
\begin{itemize}
\item time-shaped programs (4--8 segments) gained \(\le0.0025\) in \(Q\);
\item periodic schedules under continuous DNP gained \(\le0.0016\);
\item the search has a wall-time budget; when it runs out, the best validated horizon is returned
with a warning.
\end{itemize}

%==================================================================================================
\section{Software architecture and code-to-physics mapping}
\label{sec:architecture}

Paths are under \texttt{} unless they start at the repository root (marked /).

\begin{center}\small
\begin{tabularx}{0.96\textwidth}{>{\ttfamily}L{0.33\textwidth} X}
\toprule
File & Role\\
\midrule
/run\_app.py & GUI entry point; takes a material name or file, \texttt{--program}, \texttt{--state}, \texttt{--list}.\\
/design\_tensor.py & Headless tensor design from a material or snapshot.\\
/generate\_training\_data.py & Command-line dataset generator (\texttt{generate}, \texttt{info}, \texttt{export}).\\
lineshape.py & Boltzmann relations, \(P,Q\leftrightarrow\) level fractions, analytic single-site Pake branch.\\
multisite\_lineshape.py & Azimuth-averaged complex powder branch; C--D/O--D mixture.\\
model.py & \texttt{Spin1Params}, \texttt{Spin1Model}: grid, capacities, observables, mechanism profiles, pair kernels, full right-hand side.\\
spectral\_rf.py & \texttt{RFRateOperator}: \(U\mapsto W^\pm\) and its transpose.\\
ideal\_model.py & \texttt{IdealBinModel}: per-bin programs, event-aware adaptive stepping, extended-engine switch.\\
extended\_core.py & Site-resolved and dipolar-reservoir dynamics with exact VJP.\\
afp.py & AFP swap and region selection.\\
pulse\_program.py & Bin-addressed schedules, JSON/CSV I/O, compilation to an event list.\\
material\_config.py & Versioned material files and fingerprinted population snapshots.\\
presets.py & Shipped material names and paths.\\
tensor\_optimizer.py & \texttt{FrozenDynamics}, adjoint, \texttt{optimize\_tensor}, replay and validation.\\
datagen.py & Random, sweep and chain dataset generation; shard I/O.\\
playback.py & RF gate and program-start orchestration, independent of Qt.\\
baseline\_gui.py,\newline gui.py,\newline designer\_gui.py & Qt windows (layered by subclassing).\\
rf\_response\_gui.py & RF saturation / spin-diffusion parameter dialog.\\
profile\_editor.py,\newline profile\_editing.py & Per-bin program table editor.\\
/tools/steady\_state\_rf\_dnp.py & Steady-state RF\,+\,DNP optimizer.\\
/tools/plot\_tensor\_design.py & Plots of a design result.\\
/ml/ & PyTorch pulse-design and polarization models (\cref{sec:ml}).\\
/materials/, programs/ & Shipped materials and reference programs.\\
/tests/ & 460 regression tests (\cref{sec:tests}).\\
\bottomrule
\end{tabularx}
\end{center}

\subsection{Core function map}

\begin{longtable}{>{\ttfamily}L{0.36\textwidth} L{0.58\textwidth}}
\caption{Key functions and their mathematical role.}\label{tab:functionmap}\\
\toprule
Function & Meaning\\
\midrule
\endfirsthead
\toprule
Function & Meaning\\
\midrule
\endhead
normalized\_component & Normalized analytic single-site branch, \cref{eq:analyticpake}.\\
powder\_branch, site\_components & Resolvent \cref{eq:resolvent} and the two-site density \cref{eq:multisite}.\\
level\_populations\_from\_PQ & \cref{eq:pplus,eq:pzero,eq:pminus}, with a check of the physical triangle.\\
Spin1Model.reset & Grid, \(\mu_k\), initial state, reference copy, display calibration.\\
polarizations, branch\_areas & \(P\), \(Q\), \(Q_{\rm B}(P)\); branch areas.\\
spectrum\_from\_state & Projection to \(\Iplus(R)\), \(\Iminus(R)\) and total, \cref{eq:displayIplus,eq:displayIminus}.\\
branch\_indices & Physical \(R\) to packets \(k_+\) and \(k_-\), \cref{eq:kplus,eq:kminus}.\\
packet\_site\_geometry & O--D share and per-site \(\cos^2\theta\) of each packet.\\
mechanism\_profile & \(f^X_k\), \cref{eq:profile}.\\
rf\_packet\_response & \(c_k\), \cref{eq:ck}.\\
rf\_rate\_operator & Builds \code{RFRateOperator} with \(\kappa\), \(\Gamma_k\), \(c_k\).\\
capacity\_rate\_weights, dnp\_rate\_weights & DNP density weight \(w^{\rm D}_k\), \cref{eq:dnpterm}.\\
\_spectral\_overlap & \(S(\delta)\), \cref{eq:overlap}.\\
\_ensure\_diffusion\_kernel & Builds and caches the Toeplitz/Hankel FFT kernels or sparse pair lists.\\
\_spin\_diffusion\_terms & Same-branch, cross-branch and DQ currents, \cref{eq:Jsp,eq:Jsm,eq:Jx,eq:Jdq}.\\
derivative & Full right-hand side; optional breakdown by mechanism (\cref{app:breakdown}).\\
IdealBinModel.step & Event-aware adaptive Euler, \cref{eq:euler}; positivity check; exposure accounting.\\
\_max\_outgoing\_bound & \(\lambda_{\max}\) for step selection.\\
set\_program, start\_program, stop\_program & Install, anchor and stop a pulse program.\\
afp\_sweep, afp\_target\_state & AFP swap (\cref{sec:afp}), per site in the extended engine.\\
ExtendedCore.rhs / vjp & Extended dynamics and its exact adjoint.\\
PopulationSnapshot.capture / make\_model & Fingerprinted state freeze and validated restore.\\
FrozenDynamics.rhs / vjp / evaluate & Regrouped forward map, VJP, and fixed-step forward and adjoint integration.\\
optimize\_tensor & Duration scan, inner solves, selection, validation, report.\\
replay & Independent prediction through the unchanged scheduler.\\
synchronized\_program & \((U,T)\) to a \code{PulseProgram}.\\
datagen.make\_chain, undershoot & Closed-loop data generation (\cref{sec:datagen}).\\
\bottomrule
\end{longtable}

\subsection{GUI update loop}

The real-time loop is
\begin{lstlisting}[language=Python,caption={Simplified GUI update logic.}]
every 35 ms (QTimer), if not paused:
    model.step(n_steps=steps_per_tick)    # adaptive, event-split, exact pulse edges
    record I+(M,t), I-(M,t), I+(-M,t), I-(-M,t)
    if program finished and "pause at end": pause
update spectrum (+ RF overlay), tensor spectrum, monitor traces, readouts
\end{lstlisting}
The timer decides only when to draw. Pulse durations are simulation time, independent of the
redraw rate. \code{playback.py} keeps two RF conditions separate: an enabled RF gate and a
started program. Loading a file arms neither.

%==================================================================================================
\section{Graphical interface}
\label{sec:gui}

\code{run\_app.py} opens the main window, \emph{ssRF --- spin-1 tensor-polarization simulator},
paused with RF off, at the chosen material (default \code{d\_butanol\_matched}).

\subsection{Three dynamic plots}

\begin{enumerate}
  \item \textbf{Spectrum.} \(\Iplus(R)\), \(\Iminus(R)\) and their sum, with the monitor marker
  \(M\) and its mirror \(-M\). An optional overlay shows the RF commands (\code{Show RF commands
  on spectrum}).
  \item \textbf{Tensor spectrum.} \(\Iplus(R)-\Iminus(R)\). Its integral is \(Q\), and its
  negative regions are where RF raises \(Q\) (\cref{sec:rf-PQ}).
  \item \textbf{Monitor traces.} \(\Iplus\) and \(\Iminus\) at \(M\) (solid) and at \(-M\)
  (dotted), against time since the trace was started.
\end{enumerate}

\subsection{Live readouts}

\(P(t)\) and \(Q(t)\), \(Q_{\rm B}[P(t)]\), simulation time, RF/DNP/diffusion states, program
state and elapsed time, delivered exposure, and the material name and RF-response summary.

\subsection{Controls, defaults, and ranges}

Defaults are those of \code{d\_butanol\_matched}, the material the GUI opens with. The
designer window widens any range to include a value loaded from a material file. For example,
the zero-quantum width 0.717 exceeds the nominal maximum of 0.5.

\begin{longtable}{@{}L{0.14\textwidth} L{0.20\textwidth} L{0.11\textwidth} L{0.14\textwidth} L{0.30\textwidth}@{}}
\caption{GUI controls, defaults (\code{d\_butanol\_matched}), ranges, and meaning.}\label{tab:guicontrols}\\
\toprule
Group & Control & Default & GUI range & Meaning\\
\midrule
\endfirsthead
\toprule
Group & Control & Default & GUI range & Meaning\\
\midrule
\endhead
Material & load / save / save as & matched & files & Whole material (parameters, optimizer, GUI settings).\\
Material & edit ALL settings & --- & dialog & Every \code{Spin1Params} and optimizer field, validated.\\
Material & RF saturation / spin diffusion & --- & dialog & \(\kappa\), \(\Gamma\), selectivity band, exchange parameters (\cref{sec:profiles,sec:paircurrents}).\\
Material & save / load state & --- & files & Fingerprinted population snapshot (\cref{sec:snapshots}).\\
Material & generate tensor program & --- & dialog & Captures the current state and opens the designer.\\
\midrule
Run & RF ON/OFF & off & toggle & Gate for the RF term; simulation time continues.\\
Run & pause / reset / restart traces & running, paused at start & buttons & Stop stepping; rebuild the initial state; clear the monitor traces.\\
\midrule
Markers & monitor bin \(M\) & bin nearest \(R=-0.935\) & \([0,499]\) & Packet monitored by the solid traces; markers never move RF.\\
Markers & monitor \(R\) & \(-0.935\) & \([R_{\min},R_{\max}]\), step \(\Delta R\) & Snaps to the nearest bin; mirror \(-M\) is read-only.\\
\midrule
Mode & manipulation mode & ss-RF only & ss-RF / AFP / both & Which panels are active.\\
\midrule
RF program & global RF gain & 1.0 & \([0,1000]\) & Multiplies every scheduled command.\\
RF program & bin table & empty & per bin & Power \(U\), start, duration of every bin.\\
RF program & start / stop / pause at end & --- & buttons & Anchor the program at the current time; stop while evolving.\\
RF program & show RF commands & off & toggle & Overlay of \(U\) on the spectrum.\\
\midrule
AFP & region & optimal (\(Q<0\)) & optimal / selective & Bins with \(\Iplus-\Iminus<0\), or an \([L,R]\) window.\\
AFP & \(L\), \(R\) & \(-0.98\), \(-0.72\) & \([R_{\min},R_{\max}]\) & Selective window bounds.\\
AFP & efficiency \(\varepsilon\) & 1.0 & \([0,1]\) & Swap efficiency.\\
\midrule
DNP & DNP ON/OFF & off & toggle & DNP source term.\\
DNP & \(P\) saturation & 0.476 & \([-0.99,0.99]\) & \(P_{\rm sat}\).\\
DNP & DNP build rate & 0.00230 & \([0,20]\) & \(k_{\rm D}\) [s\(^{-1}\)].\\
\midrule
Lineshape & \(\Gamma\) broadening & 0.0236 & \([0.001,0.5]\) & Line width; resets the model.\\
Lineshape & \(\eta\cos2\phi\) & 0 & \([-0.5,0.5]\) & Single-site asymmetry (unused for the two-site line).\\
Lineshape & plot scale & 0.102 & \([0.01,20]\) & Display only.\\
\midrule
DNP weighting & rate-density power \(\alpha_{\rm D}\) & 1.0 & \([0,3]\) & DNP density weighting.\\
DNP weighting & max weight clip \(c\) & 12 & \([1,100]\) & Clip of the raw weights.\\
\midrule
Spin diffusion & diffusion ON/OFF & on & toggle & All exchange channels.\\
Spin diffusion & exchange scale \(K_0\) & 0.01283 & \([0,1000]\) & Flip-flop rate scale.\\
Spin diffusion & ZQ overlap shape & Lorentzian & Lor. / Gauss. & \(S(\delta)\).\\
Spin diffusion & zero-Q width \(w\) & 0.717 & \([0.001,0.5]\)\textsuperscript{*} & ZQ overlap HWHM.\\
Spin diffusion & tensor exchange ratio & 1.0 & \([0,10]\) & \(r_{\rm x}\).\\
Spin diffusion & DQ transport ratio & 3.03 & \([0,10]\) & \(r_{\rm DQ}\).\\
Spin diffusion & orientation corr. & 0 & \([0,1]\) & \(f_\theta\).\\
Spin diffusion & corr. width & 20\(^\circ\) & \([0.1,90]\) & \(\sigma_\theta\).\\
Spin diffusion & cutoff & 0 (all tails) & \([0,100]\) & In units of \(w\).\\
Spin diffusion & MW diffusion factor & 1.0 & \([0,20]\) & Multiplies \(K_0\) while DNP is on.\\
\midrule
Initial / numeric & initial \(P\) & 0.50 & \([-0.99,0.99]\) & Boltzmann start; resets the model.\\
Initial / numeric & \(T_1\) rate, \(P_{\rm eq}\) & 0, 0 & \([0,10]\), \([-0.99,0.99]\) & Spin-lattice relaxation.\\
Initial / numeric & dt & 0.03\,s & \([10^{-5},0.05]\) & Material step; substeps are automatic.\\
Initial / numeric & steps/tick & 12 & \([1,1000]\) & Material steps per 35\,ms timer tick (0.36\,s of simulation per tick).\\
Initial / numeric & plot noise & 0 & \([0,0.02]\) & Display-only noise.\\
\bottomrule
\end{longtable}
\noindent\textsuperscript{*}Widened automatically to include the loaded value.

\subsection{Backend-only parameters}

The following are stored in the material file but have no main-window control:
\begin{itemize}
  \item the grid (\code{n\_bins}=500, \code{r\_min}/\code{r\_max}=\(\mp6.186\)) and the optional initial \code{q0};
  \item the two-site lineshape parameters and \code{line\_phi\_points}=128;
  \item all mechanism-profile parameters (\cref{eq:profile}) and the extended-engine switches (site resolution, dipolar reservoir);
  \item the legacy fields (\code{gamma\_rf}=2.0, Voigt widths, \code{rf\_density\_power}, \code{rf\_pedestal\_gain}), which are unused in spectral mode;
  \item the display settings (\code{display\_scale}=0.102, \code{plot\_divisor}=10, \code{calibration\_p}=0.5);
  \item the GUI settings: timer 35\,ms, 12 steps per tick, 4500 trace points, window \(1200\times555\).
\end{itemize}

\subsection{Tensor designer window}

\emph{Generate tensor program from CURRENT state} pauses the simulation and captures a snapshot.
It opens the designer, which has the following settings:
\begin{itemize}
\item maximum command per bin;
\item duration range and number of samples;
\item iterations per start;
\item wall-time limit;
\item maximum prediction steps;
\item total power limit and total exposure limit (each optional);
\item profile smoothing;
\item candidate rule: union, tensor deficit, positive slope, or all bins;
\item exact/advanced JSON: regions, per-bin ceilings, starts, tolerances.
\end{itemize}
The calculation runs in a worker thread; cancelling leaves the live model unchanged. Three tabs
show the progress and validation log, the predicted endpoint with the RF profile, and the bin
command table. \emph{Load generated program (do not run)} installs the program with its validated
playback step as the RF integration ceiling. Export writes JSON, CSV and the snapshot.

%==================================================================================================
\section{Validation and regression tests}
\label{sec:tests}

The suite contains 460 tests (\code{python3 -m pytest}). The GUI tests skip without a display.
They establish implementation consistency and the physics invariants, not experimental validity.

\begin{longtable}{L{0.36\textwidth} L{0.58\textwidth}}
\caption{Regression-test coverage (number of tests in parentheses).}\label{tab:tests}\\
\toprule
File & Assertions\\
\midrule
\endfirsthead
\toprule
File & Assertions\\
\midrule
\endhead
test\_tensor\_designer (48) & Regrouped RHS equals the model; analytic VJP and discrete adjoint match finite differences; designs bounded, improving and replayable; refuses damaging RF; negative \(P\); cancellation leaves the live state alone; single-bin design matches the analytic optimum exposure.\\
test\_live\_diffusion\_fft (38) & FFT, sparse and pair-sum evaluators agree; every channel conserves \(\mu_k\) and \(P\) over random grids, overlaps, cutoffs and cross ratios.\\
test\_boltzmann\_recovery (30) & Each channel stationary at Boltzmann; positive entropy production; non-Boltzmann uniform states not stationary; relaxation to the unique Boltzmann line at the remaining \(P\); DQ leaves \(Q\) unchanged; RF stays an uncancelled sink during recovery; dt convergence.\\
test\_model\_diffusion\_kernel (21) & Ideal-bin mirror relation emerges from the populations; packet mass and \(P\) conservation; cross channel conserves \(P\) and changes \(Q\); post-burn currents fill holes and reduce mirror peaks; finite RF width imposes no fixed mirror ratio.\\
test\_ideal\_profiles (31), test\_custom\_editing (32), test\_playback (17), test\_rf\_integration\_clock (11) & Program editing and compilation; exact pulse edges; RF gate and program state; integration ceiling only during RF; negative \(P\) with simultaneous mirror drive.\\
test\_spectral\_saturation (18), test\_peak\_selectivity (11), test\_material\_response (17) & Spectral kernel and its transpose; selectivity and pedestal profiles; legacy and independent modes reproduce earlier laws.\\
test\_v8\_mechanisms (10), test\_v8\_extended (9) & Mechanism profiles; one-site extended engine equals the basic engine; extended adjoint at \(\beta\ne0\); exchange conserves \(E_{\rm spin}+E_d\); Provotorov sign; per-site AFP.\\
test\_multisite\_material (17), test\_production\_materials (12) & Two-site lineshape; the shipped materials load, are fingerprinted and reproduce their reference values.\\
test\_exposure\_budget (16), test\_total\_power (12), test\_prediction\_step\_limits (6), test\_profile\_refinement (11) & Caps never exceeded; SLSQP path; step plans independent of trial powers; smoothing accepted only within tolerance.\\
test\_afp (9), test\_material\_afp (30) & Mirror-safe windows; \(Q<0\) region selection; efficiency; per-site action.\\
test\_datagen (7) & Determinism; shard/manifest/resume/export round trip; late-snapshot replay; states-only; chain rounds follow applied programs; undershoot respects caps.\\
GUI tests (47) & Display projections, callbacks and dialogs (headless where possible).\\
\bottomrule
\end{longtable}

%==================================================================================================
\section{How to calibrate the simulation software against experiment}
\label{sec:calibration}

Quantitative use requires a staged calibration, not simultaneous free fitting of every
parameter. The column ``d-butanol'' states what was done for the shipped materials.

\subsection{Step 1: static lineshape and display calibration}

With RF and DNP transients absent, fit the unmanipulated spectrum for:
\begin{itemize}
  \item the frequency-to-\(R\) conversion and centre;
  \item the two-site parameters \(\gamma\), \(\rho_{\rm OD}\), \(f_{\rm OD}\), \(\eta_{\rm CD}\), \(\eta_{\rm OD}\) (or \(\Gamma_{\rm line}\) and \(s\) for one site);
  \item baseline, phase, gain and the absolute NMR calibration \(C\).
\end{itemize}
A real Q-meter fit may also require dispersion, cable response and a complex-susceptibility
treatment \cite{ClementKeller2023,Keller2026}.
\emph{d-butanol:} fitted to 2014 spectra, held-out RMS 0.79\%.

\subsection{Step 2: initial population state}

Use calibrated branch areas for \(P_0\) and, where possible, \(Q_0\). Use \(Q_0=Q_{\rm B}(P_0)\)
only if one spin temperature is verified. Otherwise load a non-Boltzmann state as a snapshot.

\subsection{Step 3: DNP-only build-up}

With RF off, fit \(P_{\rm sat}\) from the asymptote, \(k_{\rm D}\) from the time scale, and the
DNP density weighting (\(\alpha_{\rm D}\) or \(f^{\rm dnp}\)) from horn-versus-shoulder build-up
rates. These are now DNP-only parameters, so a DNP fit no longer constrains RF or recovery.
\emph{d-butanol:} \(k_{\rm D}=2.30\times10^{-3}\)\,s\(^{-1}\), \(P_{\rm sat}=0.476\), held-out
0.79\%.

\subsection{Step 4: short, low-power RF calibration}

Use burns short enough that exchange is negligible during the pulse:
\begin{itemize}
  \item \(\kappa\) from the initial fractional depletion rate;
  \item \(\Gamma\) from the narrow-hole width;
  \item per-spin selectivity (\(s\), \(f^{\rm rf}\), \(f^{\rm width}\)) from the initial
  fractional rates at horn and pedestal;
  \item the direct/mirror ratio at small \(|R|\) (\cref{sec:mirror}).
\end{itemize}
\emph{d-butanol:} \(\kappa\) was obtained jointly with exchange from one 0.1\,W burn/refill
experiment (held-out RMSE \(3.4\times10^{-4}\)). Selectivity was not measured; \(s=0.7\) is a
choice.

\subsection{Step 5: RF-off recovery}

After a calibrated burn, switch RF and microwaves off and measure:
\begin{itemize}
  \item the initial refill slope at the burn bin and its width dependence;
  \item mirror-peak decay at \(-R\);
  \item the evolution of neighbouring bins;
  \item \(P(t)\), to confirm redistribution without repolarization;
  \item \(Q(t)\), whose relaxation toward \(Q_{\rm B}(P)\) measures the cross-branch channel.
\end{itemize}
Fit \(K_0\), \(w\), \(r_{\rm x}\) and \(r_{\rm DQ}\). A narrow burn in the \(\Iminus\)-only region
(\(R\approx1.8\)) separates the channels: same-branch transport refills only \(\Iminus\) there,
while cross-branch transport also changes \(\Iplus\) at the mirror frequency.
\emph{d-butanol:} \(K_0=0.0128\), \(w=0.717\), \(r_{\rm DQ}=3.03\) from the same burn/refill fit;
\(r_{\rm x}=1\) assumed. A narrow physical width (\(w=0.047\)) fits 23\% worse.

\subsection{Step 6: DNP-on competition}

Run DNP and RF together. The calibrated terms should predict the steady hole depth, \(P\) and
\(Q\). A gentle burn at \(R\approx+1.9\) with DNP on, watching \(Q\) rise above its DNP-only value
and approach steady state over \(\sim\)435\,s, is the most direct test of the transport
assumption. \emph{d-butanol:} not yet done.

\subsection{Saturation balance: peak versus pedestal}
\label{sec:peak-pedestal}

The motivating observation is that, in deuterated butanol, a hole can need less RF power in the
pedestal region near \(\theta\simeq0\) than at a Pake peak near \(\theta\simeq90^\circ\). This
section derives what the model says about that difference and evaluates it with the calibrated
materials. Earlier spin-1 studies provide the RF/DNP and recovery context
\cite{Keller2017,KellerCrabbDay2020,ClementKeller2023,Keller2026}.

\paragraph{Physical basis.}
Khutsishvili's magnetization equation contains distinct diffusion, RF-saturation and relaxation
terms. Its stationary solution shows the competition between saturation and restoration
\cite[Eqs.~(3) and (12)]{Khutsishvili1962}. Walder \emph{et al.}\ measured the refilling of
controlled proton-magnetization holes by spin diffusion \cite{Walder2019}. Those holes are
spatial deficits in hydrated cementitious solids, so the experiment establishes a transport
mechanism and a recovery-based measurement strategy, not a deuteron peak--pedestal dependence.
For deuterium, Eckman explains why matching transitions within a local-field width permits
first-order dipolar exchange, and why quadrupolar mismatch suppresses it
\cite[Sec.~III.G]{Eckman1982}. Suter and Ernst distinguish the transport of Zeeman and
quadrupolar order in resolved spectra \cite{SuterErnst1985}.

\paragraph{A local balance.}
Use the population difference per spin in packet \(k\),
\begin{equation}
  d_{k,+}=\frac{n_{k+}-n_{k0}}{\mu_k},
  \qquad
  d_{k,-}=\frac{n_{k0}-n_{k-}}{\mu_k}.
  \label{eq:pp-normalized-difference}
\end{equation}
For an isolated addressed transition, write either difference as \(d_k\). A locally quasi-steady
approximation is
\begin{equation}
  \dot d_k=-\lambda_{{\rm RF},k}d_k
  +\kappa_k\bigl(d_{{\rm res},k}-d_k\bigr),
  \qquad
  \lambda_{{\rm RF},k}=2W_k,
  \label{eq:pp-local-balance}
\end{equation}
where \(\kappa_k\) is the effective restoring rate and \(d_{{\rm res},k}\) the reservoir
difference. For fixed coefficients,
\begin{align}
  \frac{d_{k,\rm ss}}{d_{{\rm res},k}}
    &=\frac{\kappa_k}{\lambda_{{\rm RF},k}+\kappa_k},
    \label{eq:pp-steady-balance}\\
  d_k(t)&=d_{k,\rm ss}
    +\bigl[d_k(0)-d_{k,\rm ss}\bigr]
      \exp\!\left[-(\lambda_{{\rm RF},k}+\kappa_k)t\right].
    \label{eq:pp-finite-burn}
\end{align}
Demanding a residual fraction \(\epsilon\) gives
\begin{equation}
  \lambda_{{\rm RF},k}^{\rm required}
  =\frac{1-\epsilon}{\epsilon}\,\kappa_k,
  \qquad
  U_k^{\rm required}=\frac{1-\epsilon}{\epsilon}\,\frac{\kappa_k}{\lambda_{{\rm RF},k}/U}.
  \label{eq:pp-required-rate}
\end{equation}
Stronger restoration requires more RF for equal fractional saturation. The packet capacity
cancels from this criterion; it still sets the absolute signal and the absorbed power. With
several reservoirs (exchange with unburned spins, DNP, lattice),
\begin{equation}
  d_{k,\rm ss}
  =\frac{\sum_a\kappa_{ka}d_{ka}^{\rm res}}
  {\lambda_{{\rm RF},k}+\sum_a\kappa_{ka}}.
  \label{eq:pp-multiple-reservoirs}
\end{equation}
With DNP off, \(T_1\) does not restore the hyperpolarized state, and a sustained burn depletes a
finite unpumped reservoir. In that case \(d_{\rm res}\) is not constant.

\paragraph{Converting a saturation rate to RF amplitude.}
With \cref{eq:Wgolden},
\(\lambda_{{\rm RF},k}=2\omega_{1}^2\Gamma_k/(\Gamma_k^2+\Delta_k^2)\equiv c'_kB_{1}^2\), and
\begin{equation}
  B_{1,k,\rm required}^{2}
  =\frac{1-\epsilon}{\epsilon}\,\frac{\kappa_k}{c'_k}.
  \label{eq:pp-required-field}
\end{equation}
A larger required field at the peak therefore constrains the ratio \(\kappa_k/c'_k\). It
identifies faster replenishment only once the per-spin RF response \(c'_k\) is controlled
separately. Here \(\Gamma_k\) is the homogeneous packet width. It is neither the Pake-envelope
width nor the ZQ width \(w\).

\paragraph{Replenishment from the pair currents.}
The production model gives \(\kappa_k\) explicitly. Take a small \(\Iplus\) hole in packet \(k\),
\(p_{k+}=\bar p_+-\delta\), \(p_{k0}=\bar p_0+\delta\), with all other packets at the Boltzmann
fractions \(\bar p\). Linearizing \cref{eq:Jsp,eq:Jsm} gives
\begin{equation}
  \kappa^{\rm SQ}_k = K_0\,\bigl[(S\mu)_k-\mu_k\bigr]\Bigl(1-\frac{\bar p_-}{2}\Bigr),
  \qquad (S\mu)_k=\sum_jS(x_k-x_j)\,\mu_j .
  \label{eq:kappaSQ}
\end{equation}
The same-branch refill rate is proportional to the \emph{number of partners within the ZQ width},
\((S\mu)_k-\mu_k\). The cross-branch rate is proportional to the partners at the opposite
frequency, \((X\mu)_k\). The DQ rate is proportional to the global partner sum, which is the same
for every packet. For a narrow ZQ width, \((S\mu)_k\) follows the local Pake density, and a horn
packet has many more partners than a pedestal packet. For a broad width it is nearly uniform. The
powder density thus enters replenishment through partner counting at a specified width. That is
a definite, testable mechanism, not an imposed weight.

\paragraph{Calibrated predictions.} \Cref{tab:peak-pedestal} evaluates \(\kappa_k\) (from the full
model, linearized) and the per-spin RF rate for a narrow \(\Iplus\) hole. The state is Boltzmann
at \(P=0.45\) with DNP off. The peak is at \(R=0.93\); the pedestal at \(R=-1.60\), on the same
branch toward \(\theta=0\).

\begin{table}[h]
\centering\footnotesize\setlength{\tabcolsep}{4pt}
\begin{tabular}{lcccccc}
\toprule
Material & \multicolumn{2}{c}{\(\kappa_k\) [s\(^{-1}\)]} & \multicolumn{2}{c}{\(\lambda_{\rm RF}/U\) [s\(^{-1}\)]} & \(U^{\rm req}\) ratio & Peak channels\\
 & peak & pedestal & peak & pedestal & peak/ped. & SQ / X / DQ\\
\midrule
matched (\(w=0.717\), \(r_{\rm DQ}=3.03\)) & 0.0114 & 0.0079 & 2.61 & 2.61 & 1.4 & 0.0061 / 0.0028 / 0.0026\\
same-branch (\(w=0.047\), SQ only) & 0.0435 & 0.0101 & 1.44 & 1.44 & 4.3 & 0.0435 / 0 / 0\\
extend (horn-favouring RF) & 0.0114 & 0.0079 & 14.2 & 0.64 & 0.065 & as matched\\
\bottomrule
\end{tabular}
\caption{Peak--pedestal saturation balance from the calibrated materials, \cref{eq:pp-required-rate}.
\(U^{\rm req}\) is the command needed for a given fractional saturation.}
\label{tab:peak-pedestal}
\end{table}

Three conclusions follow:
\begin{enumerate}
  \item With equal per-spin RF rates (the matched and same-branch materials), the peak needs
  \emph{more} RF than the pedestal. The ratio is mild (1.4) with the fitted broad, nonlocal
  transport and large (4.3) with narrow, local same-branch transport. The observed contrast is
  thus naturally explained by \emph{partner counting at a narrow ZQ width}.
  \item Horn-favouring per-spin RF selectivity (\code{d\_butanol\_extend}) predicts the
  \emph{opposite}: the peak needs about 15 times less RF. That material maximizes single-burn
  \(Q\), but it is disfavoured if the peak--pedestal observation holds.
  \item The observation therefore discriminates between two unconstrained parameters that matter
  most for tensor enhancement: the per-spin selectivity and the locality of transport. A
  quantitative peak--pedestal measurement (required power at matched fractional depth, plus
  initial refill rates) should be the next calibration experiment.
\end{enumerate}
These are linearized, single-bin rates with an unburned reservoir. A finite hole width depletes
the nearest partners too, which lowers the effective \(\kappa\). This happens most for narrow
\(w\), so the full-model required power should be computed for the actual burn width
(\code{tools/steady\_state\_rf\_dnp.py} and the GUI support this directly).

\paragraph{Controlled simulation comparisons.}
\begin{enumerate}
  \item \textbf{Compare the same transition at its peak and pedestal.} Compare \(\Iplus\) near
  \(R=+1\) with its pedestal toward \(R=-2\), or \(\Iminus\) near \(R=-1\) with its pedestal toward
  \(R=+2\). Comparing \(\Iplus\) and \(\Iminus\) at one \(R\) tests overlap selectivity instead.
  Hold duration, bandwidth and calibrated drive fixed, and compare fractional depths. For a
  required-drive comparison, fix the target depth and vary the drive.
  \item \textbf{Capacity-only baseline.} With all profiles at 1 and exchange, DNP and \(T_1\) off,
  \(d_k(t)/d_k(0)=\exp(-2W_kt)\) is the same at every location despite unequal absolute signals.
  This is the null check.
  \item \textbf{Separate RF response from transport.} Vary \(f^{\rm rf}\) or \(f^{\rm width}\)
  with exchange off, then vary \(K_0\), \(w\) and the channel ratios with RF fixed. Use the
  derivative breakdown (\cref{app:breakdown}) to record each channel's contribution.
  \item \textbf{Unpumped versus maintained reservoirs.} Compare burns and dark recovery with DNP
  off (retaining the reduced \(P\)) and with DNP on, after verifying the driven steady state.
  \item \textbf{Time, bandwidth and acquisition.} Compare early depletion, late depth and recovery
  for matched burns. Vary the burn width, and include readout delay or sweep averaging when
  comparing with data. Under grid refinement, preserve the physical bandwidth and the total
  command.
\end{enumerate}

%==================================================================================================
\section{Calibrated materials}
\label{sec:materials}

\begin{table}[h]
\centering\small
\begin{tabular}{lll}
\toprule
Quantity & \code{d\_butanol\_matched} & Source \\
\midrule
Grid & \(N=500\), \(|R|\le6.186\) & covers the O--D line \\
Line & two-site, \(\gamma=0.0236\), \(\rho_{\rm OD}=1.355\), \(f_{\rm OD}=0.021\) & 2014 spectra, held-out RMS 0.79\% \\
& \(\eta_{\rm CD}=0.065\), \(\eta_{\rm OD}=0.033\) & \\
RF & \(\kappa=0.0196\)\,s\(^{-1}\)U\(^{-1}\), \(\Gamma=0.015\) & 0.1\,W burn/refill, RMSE \(3.4\times10^{-4}\) \\
Selectivity & \(s=0.7\) at \(R_c=0.935\) (half-width 0.16, taper 0.25) & chosen, not fitted \\
Exchange & \(K_0=0.0128\), Lorentzian \(w=0.717\) & same burn/refill fit \\
& \(r_{\rm x}=1\), \(r_{\rm DQ}=3.03\) & \(r_{\rm x}\) never fitted \\
DNP & \(k_{\rm D}=2.30\times10^{-3}\)\,s\(^{-1}\), \(P_{\rm sat}=0.476\), \(\alpha_{\rm D}=1\) & 2014 build-up, held-out 0.79\% \\
Step & \(\delta t=0.03\)\,s & \\
Units & time in seconds; \(\sum U=1\) is 0.1\,W & event IDs read as seconds \\
\bottomrule
\end{tabular}
\caption{The experimental fit.}
\label{tab:materials}
\end{table}

The other shipped materials differ only where the data are silent.

\begin{longtable}{L{0.28\textwidth} L{0.40\textwidth} L{0.24\textwidth}}
\toprule
Material & Differences from \code{matched} & Best single-burn \(Q\) from \(P_0=0.50\)\\
\midrule
\code{d\_butanol\_matched} & --- & 0.245--0.248 \\
\code{d\_butanol\_basic} & \(s=1\); RF coupling density exponent 0.5; RF width \(a_2=0.6\) & 0.265 \\
\code{d\_butanol\_extend} & \(s=1\); coupling exponent 1; width density exponent \(-0.3\); \(a_2=0.6\) & 0.275 \\
\code{d\_butanol\_full\_physics} & site resolution and dipolar reservoir (\(c_d=0.01\)) & as matched \(\pm0.0003\) \\
\code{d\_butanol\_same\_}\newline\code{branch\_transport} & hypothesis: narrow width \(w=2\Gamma=0.047\); refit \(K_0=0.516\), \(\kappa=0.0108\), \(k_{\rm D}=2.28\times10^{-3}\), \(P_{\rm sat}=0.451\); \(r_{\rm x}=0\), \(r_{\rm DQ}=0\) & steady RF+DNP 0.241 \\
\code{demo\_single\_pake} & single site on \([-3,3]\); uncalibrated & --- \\
\bottomrule
\end{longtable}

\noindent The largest model levers on tensor enhancement are:
\begin{itemize}
\item per-spin RF selectivity, about 0.04 in single-burn \(Q\) (unidentified);
\item the type and strength of transport under continuous DNP, up to 0.09 (weakly identified);
\item the RF law, 0.015--0.02 (the density double counting of earlier versions).
\end{itemize}
Calibration, the power cap and the lineshape each move \(Q\) by less than 0.01. The repository
page on key results gives the full list.

%==================================================================================================
\section{Command-line training-data generation}
\label{sec:datagen}

\code{generate\_training\_data.py} turns the simulator and optimizer into a deterministic dataset
factory for machine learning. Each sample pairs an input with a target:
\begin{itemize}
\item \textbf{Input:} the total spectrum of the input state (\cref{sec:observables}), optionally
with the branch signals and level populations.
\item \textbf{Target:} the optimizer's \((U,T)\), with its predicted \(P\) and \(Q\), the no-RF
control, the validated playback step, and optionally the post-pulse evolution.
\end{itemize}

\subsection{State distributions}

\begin{itemize}
  \item \textbf{Random mode} (\code{--samples}). A Boltzmann start at random \(P_0\), then a random
  history: log-uniform burn windows, powers and durations; optional AFP; free evolution with DNP
  on, off or random.
  \item \textbf{Sweep mode} (\code{--p0-values}). A grid of \(P_0\) \(\times\) random burns
  \(\times\) fixed evolution times. All times of one burn share its history, and the time since
  the burn is a model input.
  \item \textbf{Chain mode} (\code{--chains}, closed loop). Optimize from the current state (one
  sample). Apply an \emph{undershot} version of that program. Evolve. Repeat for \code{--rounds}
  rounds.
\end{itemize}

The undershoot in chain mode works as follows:
\begin{itemize}
\item the power is scaled by \(a\sim\mathcal U[0.3,1]\) and the duration by
\(b\sim\mathcal U[0.2,1]\);
\item the per-bin profile gets log-normal noise (\(\sigma=0.5\)), is renormalized, and is clipped
to the program's own maximum;
\item with probability 0.15 the exact optimum is applied instead;
\item the evolution time is log-uniform in 2--40\,s, the NMR and control-loop latency.
\end{itemize}
Undershooting only lowers power and time, so every applied pulse satisfies the caps. The time
since the last pulse is recorded as a model input.

\subsection{Why closed-loop data}

In a 240-sample sweep pilot (\(P_0\) 25--35\%, DNP on, 0.3\,W cap), the optimal profile barely
depended on the state:
\begin{itemize}
\item the cosine similarity of every target to the mean was at least 0.982;
\item a single fixed program reached 99\% of the optimizer's gain.
\end{itemize}
In chain data the optimum changes from round to round:

\begin{center}
\begin{tabular}{lcccc}
\toprule
Round & Q before pulse & Q with optimum & Mean duration & Power below \(R=0\)\\
\midrule
0 (fresh) & 0.068 & 0.125 & 60\,s (bound) & 29\%\\
1 & 0.114 & 0.130 & 28\,s & 43\%\\
2 & 0.119 & 0.132 & 20\,s & 46\%\\
3 & 0.124 & 0.133 & 14\,s & 46\%\\
\bottomrule
\end{tabular}
\end{center}

\noindent Re-applying the pilot's mean program on later-round states gave 75\% of the gain on
average, and sometimes ended with lower \(Q\) than no RF.

\subsection{Implementation}

\begin{itemize}
\item \textbf{Determinism.} Every random draw comes from
\code{np.random.default\_rng([seed, tag, unit, k])}, so any sample can be regenerated alone.
\item \textbf{Shards.} Work units (samples, sweep groups or chains) are written as \code{.npz}
shards through a temporary file and an atomic rename, so a shard is either complete or absent.
\item \textbf{Manifest and resume.} A manifest stores the configuration, the full material files
and the grid. A run resumes only if its configuration is identical.
\item \textbf{Parallelism.} Workers are spawned processes with one BLAS thread each.
\item \textbf{Optimizer presets} bound the wall time. The \code{fast} preset uses 5 durations, 1
refinement, 40 iterations and 1 start; at the 0.3\,W cap one design takes 1--3 CPU-minutes, about
200 samples per hour on 7 workers.
\item \textbf{States only} (\code{--states-only}) skips the optimizer and stores spectra with
exact \(P\), \(Q\) and branch labels, for polarization models, at about 0.5\,CPU-s per sample.
\end{itemize}

%==================================================================================================
\section{Learning models}
\label{sec:ml}

\subsection{Pulse-design model}

\code{SpectrumToPulse} is a 1D U-Net over the frequency axis:
\begin{itemize}
\item \textbf{Encoder:} three down-sampling stages to \(N/8\); convolutions with kernel 5,
GroupNorm and GELU.
\item \textbf{Decoder:} skip connections; a \(1\times1\) convolution gives profile logits \(\ell\).
\item \textbf{Scalar head:} an MLP on the mean- and max-pooled bottleneck gives the total power
and \(\log T\).
\item \textbf{Inputs:}
  \begin{itemize}
  \item the spectrum with one \emph{global} scale, since the absolute amplitude carries \(P\);
  \item its mirror \(R\to-R\), since each frequency couples to its mirror through the packets;
  \item in sweep and chain data, a constant channel with the time since the last pulse.
  \end{itemize}
\item \textbf{Output:} \(U=\bigl(\sum U\bigr)\,\mathrm{softmax}(\ell)\), applied for \(T\).
\item \textbf{Loss:} KL divergence on the profile, plus \(0.5\times\) MSE on the standardized
total and \(\log T\).
\end{itemize}
With \code{--residual}, \(\ell=\log\bar\pi+\ell_\theta\), where \(\bar\pi\) is the mean training
profile and the last layer is zero-initialized. The untrained network is then exactly the mean
program, and training learns only the state-dependent correction.

\paragraph{Physics evaluation.} Each prediction is first projected onto the caps: rescale
\(\sum U\) to the power cap, clip \(T\), and shorten \(T\) to the energy cap. The projected
program is then simulated from the stored populations. The score is the gain fraction
\((Q_{\rm net}-Q_{\rm noRF})/(Q_{\rm opt}-Q_{\rm noRF})\), compared with two baselines:
\begin{itemize}
\item the mean training program;
\item a nearest-neighbour lookup of the training spectrum.
\end{itemize}

\begin{center}
{\small\setlength{\tabcolsep}{4pt}
\begin{tabular}{llccc}
\toprule
Training data & Test & Network & Nearest neighb. & Mean program\\
\midrule
Sweep, 240 & held-out \(P_0\) 29, 33\% (plain) & 0.892 & 0.999 & 0.991\\
Sweep, 240 & same, \code{--residual} & 0.997 & 0.999 & 0.991\\
Chain, 70\(\times\)6 & 14 new chains, residual & \textbf{0.991} (worst 0.90) & 0.998 & 0.926 (worst 0.53)\\
Chain, 70\(\times\)6 & same, plain & 0.978 (worst 0.52) & 0.998 & 0.926\\
\bottomrule
\end{tabular}}
\end{center}

\noindent On closed-loop states the network learns state-specific pulses. Nearest-neighbour
lookup is still slightly better at this dataset size, and it is the benchmark to beat as datasets
grow.

\subsection{Polarization model}

\code{SpectrumToPQ} uses the same trunk. A pooled head outputs \((P,Q)\), and a per-bin auxiliary
head outputs \(\Iplus\), with \(\Iminus=I-\Iplus\). Supervising the branch split encodes
\(Q=A_+-A_-\) (\cref{eq:qarea}). The baselines are the area method for \(P\) and \(Q_{\rm B}(P)\)
of \cref{eq:QB}. The Boltzmann relation is exact for unmanipulated states, so a useful network
must beat it on manipulated ones.

%==================================================================================================
\section{Settings used in the present simulations}
\label{sec:settings}

This section records the values behind every result in this note, so that each can be
reproduced exactly.

\subsection{Full parameter set of \texttt{d\_butanol\_matched}}

\begin{longtable}{>{\ttfamily}L{0.36\textwidth} L{0.24\textwidth} L{0.32\textwidth}}
\caption{Every kinetic parameter of the default material.}\label{tab:fullparams}\\
\toprule
Parameter & Value & Note\\
\midrule
\endfirsthead
\toprule
Parameter & Value & Note\\
\midrule
\endhead
n\_bins, r\_min, r\_max & 500, \(\mp6.18620\) & \(\Delta R=0.024794\)\\
line\_model & multisite\_butanol & \\
line\_gamma & 0.0235706 & \(\gamma\)\\
line\_cd\_scale\_R & 1.0 & C--D unit of \(R\)\\
line\_od\_scale\_ratio & 1.354745 & \(\rho_{\rm OD}\)\\
line\_eta\_cd, line\_eta\_od & 0.0647698, 0.0326686 & \\
line\_od\_fraction & 0.0212736 & \(f_{\rm OD}\)\\
line\_phi\_points & 128 & azimuthal quadrature\\
p0, q0 & 0.50, Boltzmann & initial state\\
rf\_response\_mode & spectral & \\
rf\_spectral\_strength & 0.0195741 & \(\kappa\) [R\,s\(^{-1}\)U\(^{-1}\)]\\
rf\_homogeneous\_hwhm\_R & 0.015 & \(\Gamma\)\\
rf\_peak\_selectivity & 0.7 & \(s\)\\
rf\_peak\_center\_R & 0.935230 & \(R_c\)\\
rf\_peak\_half\_width\_R, rf\_transition\_width\_R & 0.16, 0.25 & selectivity band\\
rf\_coupling\_*, rf\_width\_*, exchange\_*, dnp\_*, t1\_* profiles & neutral (exponents 0, \(a_{2,4}=0\), O--D factor 1, clip 12) & \(f^X\equiv1\)\\
diffusion\_enabled & true & \\
diffusion\_scale & 0.0128319 & \(K_0\) [s\(^{-1}\)]\\
zq\_width\_R, diffusion\_overlap & 0.716884, lorentzian & \(w\), \(S\)\\
cross\_branch\_ratio & 1.0 & \(r_{\rm x}\)\\
double\_quantum\_ratio & 3.032153 & \(r_{\rm DQ}\)\\
orientation\_corr\_fraction, \_width\_deg & 0, 20 & off\\
kernel\_cutoff\_widths & 0 & full tails\\
microwave\_diffusion\_factor & 1.0 & \\
p\_dnp\_sat, dnp\_rate & 0.475633, 0.00229831 & \(P_{\rm sat}\), \(k_{\rm D}\) [s\(^{-1}\)]\\
capacity\_rate\_power, \_clip & 1.0, 12 & DNP weighting\\
t1\_rate, t1\_p\_eq & 0, 0 & off\\
site\_resolved, dipolar\_enabled & false, false & basic engine\\
afp\_efficiency & 1.0 & \\
dt & 0.03 & s\\
\bottomrule
\end{longtable}

\noindent The other materials change only the following:
\begin{itemize}
\item \code{d\_butanol\_basic}: \code{rf\_peak\_selectivity} 1.0, \code{rf\_coupling\_density\_exponent} 0.5, \code{rf\_width\_theta\_p2} 0.6.
\item \code{d\_butanol\_extend}: selectivity 1.0, coupling exponent 1.0, \code{rf\_width\_density\_exponent} \(-0.3\), \code{rf\_width\_theta\_p2} 0.6.
\item \code{d\_butanol\_full\_physics}: \code{site\_resolved} and \code{dipolar\_enabled} true, \code{dipolar\_capacity\_ratio} 0.01 (\(\xi=1\), \(r_{1d}=0\), \(\beta_L=\beta_0=0\)).
\item \code{d\_butanol\_same\_branch\_transport}: \(\kappa=0.0108192\), \(K_0=0.515827\), \(w=0.0471\), \(r_{\rm x}=r_{\rm DQ}=0\), \(P_{\rm sat}=0.450859\), \(k_{\rm D}=0.00227818\); monitor at \(R=1.9\).
\item \code{demo\_single\_pake}: single site on \([-3,3]\), \(\Gamma_{\rm line}=0.05\), \(s=0.04\), \(\kappa=0.04\), \(R_c=0.96\), \(K_0=5\), \(w=0.05\), \(r_{\rm DQ}=0.1\), \(P_{\rm sat}=0.58\), \(k_{\rm D}=0.05\); arbitrary units.
\end{itemize}

\subsection{Optimizer settings}

\begin{longtable}{>{\ttfamily}L{0.29\textwidth} L{0.15\textwidth} L{0.15\textwidth} L{0.29\textwidth}}
\caption{Optimizer settings: material default (GUI designer) and the \code{fast} preset with the
pilot configuration (data generation).}\label{tab:optsettings}\\
\toprule
Setting & Material default & Data generation & Meaning\\
\midrule
\endfirsthead
\toprule
Setting & Material default & Data generation & Meaning\\
\midrule
\endhead
max\_power & 20 & 20 & per-bin ceiling \(\bar U\)\\
max\_total\_power & none & 3 (0.3\,W) & \(P_{\rm cap}\)\\
max\_total\_exposure & none & 102 (10.2\,J) & \(E_{\rm cap}\)\\
min\_duration, max\_duration & 0.01, 10\,s & 0.5, 60\,s & duration range\\
duration\_samples & 9 & 5 & geometric scan\\
duration\_refinements & 3 & 1 & bisections\\
max\_iterations & 100 & 40 & per start\\
starts & 2 & 1 & seeds\\
search\_dt & 0.004 & 0.03 & search step ceiling\\
max\_steps & 20000 & 20000 & per horizon\\
max\_validation\_halvings & 12 & 12 & playback refinement\\
max\_wall\_seconds & 300 & 180 & time budget\\
q\_absolute / relative\_tolerance & \(10^{-6}\), 0.002 & same & selection tolerance\\
convergence\_tolerance & \(2\times10^{-5}\) & same & time-step validation\\
candidate\_mode & union & union & candidate rule\\
profile\_smoothing\_fwhm\_R & 0.03 & 0 & smoothing\\
seed & 42 & per sample & \\
\bottomrule
\end{longtable}

\noindent The \code{standard} and \code{full} presets keep the material settings, with wall budgets
of 600\,s and 1800\,s; \code{test} is for unit tests only.

\paragraph{Study conventions for the key results.}
Single-burn results start from a Boltzmann state at \(P_0=0.50\) (\(Q_0=0.1972\)) with DNP off.
They use two RF limits:
\begin{itemize}
\item \emph{uncapped}: \(U_j\le20\), with no total power or exposure limit;
\item \emph{capped}: \(\sum U\le1\) (0.1\,W) and \(T\sum U\le34\) (3.4\,J).
\end{itemize}
Continuous-DNP results set the DNP target to \(P=0.50\), so DNP alone holds \(Q=0.1972\).

\subsection{Dataset generation runs}

All three runs used \code{d\_butanol\_matched}, \code{--dnp on} (DNP during history, waits,
optimized pulse and post-pulse evolution), \code{--power-cap 3}, \code{--exposure-cap 102},
\code{--preset fast} and \code{--config data/pilot\_config.json}:
\begin{lstlisting}
{"design": {"overrides": {"min_duration": 0.5, "max_duration": 60.0,
                          "search_dt": 0.03, "max_steps": 20000}},
 "history": {"wait": [0.0, 60.0]}}
\end{lstlisting}

\begin{longtable}{L{0.17\textwidth} L{0.52\textwidth} L{0.23\textwidth}}
\toprule
Run & Settings & Size and cost\\
\midrule
Sweep pilot (\code{sweep\_dnp\_on}) & \code{--p0-values 25 27 29 31 33 35}, \code{--histories-per-p0 8}, \code{--times 0 2 5 10 20}; one burn per history (centre \(|R|\in[0,2.2]\), width 0.02--0.6, total power 0.01--5\,U, duration 0.01--20\,s, all log-uniform); post-pulse 30\,s in 3 steps with DNP; seed 12345 & 240 samples, all completed; 8714\,s on 7 workers\\
Chain probe (\code{chain\_probe}) & \code{--chains 14 --rounds 4 --p0-range 25 35}; history 0--1 burns (no wait) and, with probability 0.2, one AFP window (width 0.05--0.6, efficiency 0.7--1); no post-pulse recording; round wait 2--40\,s log-uniform; undershoot power 0.3--1, duration 0.2--1, profile noise 0.5, exact 15\%; seed 12345 & 56 samples, all completed; 1015\,s\\
Chain dataset (\code{chain\_v1}) & as the probe, \code{--chains 70 --rounds 6}, seed 2026 & 420 samples, all completed; 7262\,s (\(\approx\)208/h); mean design 117\,s\\
\bottomrule
\end{longtable}

\subsection{Training runs}

All runs used \code{python3 -m ml.train}: U-Net width 32, batch 32, AdamW (weight decay
\(10^{-4}\)) with cosine decay, loss weight 0.5 on the scalars, the best validation checkpoint,
seed 0, the Apple GPU (\code{mps}). The physics evaluation integrates at \(\min(0.03\,\mathrm{s},\,0.2/\lambda_{\max})\).

\begin{longtable}{L{0.27\textwidth} L{0.42\textwidth} L{0.23\textwidth}}
\toprule
Run & Options & Split\\
\midrule
\code{ml\_runs/sweep} & \code{--data sweep\_dnp\_on --epochs 300} (lr \(2\times10^{-3}\)) & by group 80/10/10\\
\code{ml\_runs/sweep\_holdout} & \code{--holdout-p0 29 33 --lr 5e-4 --epochs 300} & test = \(P_0\) 29, 33\%\\
\code{ml\_runs/sweep\_}\newline\code{holdout\_residual} & as above, \code{--residual} & same\\
\code{ml\_runs/chain\_}\newline\code{v1\_residual} & \code{--data chain\_v1} \code{--test-data chain\_probe} \code{--epochs 400} \code{--lr 5e-4} \code{--residual} & train/val 372/48; test 56\\
\code{ml\_runs/chain\_v1\_plain} & as above without \code{--residual} & same\\
\bottomrule
\end{longtable}

\subsection{Peak--pedestal evaluation}

\Cref{tab:peak-pedestal} was computed in the following setting:
\begin{itemize}
\item the full model's derivative, with DNP and RF off;
\item a Boltzmann state at \(P=0.45\);
\item a relative perturbation \(10^{-4}\) of the \(a\)-mode in the packet nearest \(R=0.935\) (peak) or \(R=-1.6\) (pedestal);
\item \(\lambda_{\rm RF}/U=2W^+_k\) for a unit command on that bin.
\end{itemize}

%==================================================================================================
\section{Exact physics, phenomenological assumptions, and limitations}
\label{sec:limits}

\subsection{Relations that are exact within the spin-1 packet model}

\begin{itemize}
  \item The definitions of \(P\), \(Q\) and the three populations, and the transition intensities as population differences.
  \item The mirror relation between the two spectral appearances of a packet.
  \item Conservation of packet spin count by every term.
  \item The equality of the two single-quantum RF matrix elements, and the one-half mirror rule for an isolated narrow carrier.
  \item \(\dot P|_{\rm RF}=-\sum_k(J^+_k+J^-_k)\) and \(\dot Q|_{\rm RF}=3\sum_k(J^-_k-J^+_k)\).
  \item The channel structure of spin-1 flip-flops, their equal matrix elements, and their conservation laws.
  \item The H-theorem and the uniqueness of the Boltzmann stationary state at the remaining \(P\).
  \item Energy conservation of exchange with the dipolar reservoir.
  \item The discrete adjoint gradient of the optimizer's forward map.
\end{itemize}

\subsection{Phenomenological choices}

\begin{itemize}
  \item The ZQ overlap shape and width, and the effective DQ ratio (a second-order process folded into one constant).
  \item Spatial mean-field pairing: the frequency-space kernel carries no real-space structure unless orientation correlation is enabled.
  \item The per-spin RF coupling and width profiles; the peak-selectivity band.
  \item DNP as relaxation to a Boltzmann target, with an empirical density weighting.
  \item The first-order (in \(\beta\delta\)) dipolar reservoir.
  \item Instantaneous AFP.
  \item One effective \(T_1\), with an optional profile.
\end{itemize}

\subsection{Important missing physics and extensions}

\begin{enumerate}
  \item \textbf{Coherent RF.} Nutation, spin locking and finite-rate passage, for \(\omega_1\gtrsim\Gamma\).
  \item \textbf{Real-space transport.} Diffusion barriers around paramagnetic centres; spatially inhomogeneous DNP and microwave field.
  \item \textbf{Explicit DNP channels.} Solid effect, cross effect and thermal mixing as separate electron--nuclear rate terms, and their combination when several act at once.
  \item \textbf{Exponential reservoir coupling} for large \(\beta\delta\).
  \item \textbf{Complex CW-NMR response.} Absorption, dispersion, Q-meter phase, sweep timing and acquisition latency, for direct comparison with raw data.
  \item \textbf{Global or time-shaped optimization.} The present designer is a local search over synchronized programs.
  \item \textbf{The persistent \(\sim\)20\% RF-on plateau at the peak} seen in data is not reproduced by any present parameter set.
\end{enumerate}

New mechanisms enter at three places:
\begin{enumerate}
\item a term in \code{Spin1Model.derivative} (or \code{ExtendedCore.rhs});
\item the same term, regrouped, in \code{FrozenDynamics.rhs}, together with its VJP;
\item its contribution to the outgoing-rate bound.
\end{enumerate}
The existing tests are the template: compare the evaluators, check the adjoint against finite
differences, and check the conservation laws.

%==================================================================================================
\section{Changes since the first release}
\label{sec:changes}

\begin{longtable}{L{0.18\textwidth} L{0.36\textwidth} L{0.38\textwidth}}
\toprule
Aspect & First release & Production release \\
\midrule
RF drive & one ideal bin at one burn location; rate \(\Gamma_{\rm RF}w_k\) & arbitrary per-bin, time-scheduled programs; Lorentzian homogeneous kernel; per-spin coupling and width profiles \\
Density weighting & \(w_k\propto\rho_k^\alpha\) on RF, DNP and recovery, on top of \(\mu_k\) & capacity \(\mu_k\) only, for RF and exchange; explicit mechanism profiles; density weight kept for DNP \\
Recovery & linear relaxation of two RF modes toward \(\bm n^{\rm B}[P(t)]\), plus nearest-neighbour Gaussian diffusion & mass-action pair currents for same-branch, cross-branch and DQ channels on a nonlocal Lorentzian kernel \\
Conservation of \(P\) & enforced by projecting each internal term & by construction, pair by pair \\
Approach to Boltzmann & imposed by the reference state & derived: the unique stationary state \\
Sites, reservoir & none & optional C--D/O--D populations and dipolar reservoir \\
AFP & none & instantaneous swap with efficiency \\
Integration & fixed step; silent positivity clip & rate-bounded substeps, event splitting, error on real negatives \\
Design & none & adjoint-based, cap-constrained tensor optimizer with replay validation \\
Data / ML & none & command-line generator (random, sweep, chain) and PyTorch models \\
Calibration & dimensionless & d-butanol in seconds and watts \\
Grid & \(N=701\) on \([-3,3]\) & \(N=500\), symmetric, on \([-6.19,6.19]\) for d-butanol \\
Tests & 13 & 460 \\
\bottomrule
\end{longtable}

Two of these changes alter the physics, not only the numerics.

\emph{First}, density is no longer counted twice. In the first release, \(w_k\) multiplied rates
that already scale with \(\mu_k\) through the populations, so absolute horn currents scaled
roughly as \(\rho^2\). The same double counting, in the RF law of the MaterialTuning-V1
development version, accounts for most of the difference between that version and the current
per-spin law: \(Q\approx0.264\) versus 0.245--0.248 from \(P_0=0.5\). The density law is still
available, but as an explicit hypothesis about per-spin coupling.

\emph{Second}, the first-release recovery relaxed toward a Boltzmann reference at the current
\(P\), which pulled \(Q\) toward \(Q_{\rm B}(P)\) by construction. Its separate-mode neighbour
flux also conserved \(Q\) exactly. In the production model, tensor order at fixed \(P\) relaxes
only through the cross-branch channel, so its rate is a material parameter that can be measured.
Whether real transport is same-branch or cross-branch is the largest open question for steady
operation (\cref{sec:regimes}) and for the peak--pedestal contrast (\cref{sec:peak-pedestal}).

%==================================================================================================
\section{Conclusions}

The ss-RF package models a spin-1 Pake doublet at the level of packet populations. Two
structural features carry over from the first release:
\begin{itemize}
\item the two direct components at one frequency are different orientation packets;
\item every direct hole is tied to its mirror response through shared populations.
\end{itemize}

The production release replaces the phenomenological recovery with the three spin-1 flip-flop
channels, written as reversible pair currents. As a result, conservation of vector polarization,
the distinction between spectral recovery and repolarization, and the approach to the Boltzmann
line at the remaining \(P\) are consequences of the model rather than constraints imposed on it.
The capacity of a packet and the rate per spin are kept separate: every per-spin dependence on
frequency, orientation or site is an explicit, separately tunable profile.

On this model the software builds:
\begin{itemize}
\item a validated, adjoint-based designer of \(Q\)-maximizing RF programs under realistic power
and energy limits;
\item a closed-loop data generator that reproduces the operating cycle;
\item learned models that reach 99\% of the optimizer's gain on independent closed-loop states.
\end{itemize}

The calibrated d-butanol materials fit the lineshape, DNP build-up and one burn/refill
experiment. Two quantities that matter most for tensor enhancement remain unconstrained: the
per-spin RF selectivity and the locality of spin-diffusion transport. The peak--pedestal analysis
shows that one experiment separates them. It should be done next: the required RF power at a
matched fractional hole depth at peak and pedestal, together with the initial refill rates and a
refill measurement in the \(\Iminus\)-only region.

\appendix

%==================================================================================================
\section{Population-level differential rate equations and observable intensity evolution}
\label{app:effective-rate-equations}

This appendix collects the equations used by the simulation software and derives the quantities
displayed in the graphical front end. It begins with the three spin-1 level populations in each
orientation packet. It keeps the two packets that overlap at one physical frequency separate,
and ends with the time derivatives of the plotted absorption components \(\Iplus(R,t)\) and
\(\Iminus(R,t)\). The RF bookkeeping and the direct-hole/mirror-peak relation follow
Ref.~\cite{ClementKeller2023}. The recovery equations implement the spectral-overlap transport
of Ref.~\cite{Keller2026} as pair currents.

\subsection{The state variable: three populations in each spin packet}

A spin-1 nucleus has three Zeeman sublevels \(m=+1,0,-1\). The simulation divides the broadened
line into \(N\) packets indexed by \(k\), at \(x_k=R_+(\theta_k)\), each with the population
vector
\begin{equation}
  \bm n_k(t)
  =
  \begin{pmatrix}
    n_{k+}(t)\\ n_{k0}(t)\\ n_{k-}(t)
  \end{pmatrix},
  \qquad
  n_{k+}+n_{k0}+n_{k-}=\mu_k,
  \qquad
  \sum_{k=0}^{N-1}\mu_k=1.
  \label{eq:app-packet-normalization}
\end{equation}
The fractional populations within a packet are \(p_{km}=n_{km}/\mu_k\). The \(n_{km}\) are
absolute population fractions of the whole sample; the \(p_{km}\) describe how the spins inside
packet \(k\) are divided among the levels. With the site-resolved engine, the state is
\(n_{s,k,m}\) with \(\sum_sn_{s,k,m}=n_{km}\).

\subsection{Why one physical frequency contains two different population vectors}

Packet \(k\) appears at \(R=x_k\) in \(\Iplus\) and at \(R=-x_k\) in \(\Iminus\). A frequency
\(R_b\) in the overlap region therefore contains two generally different packets:
\begin{equation}
  A\equiv k_{+}(R_b):\ x_A\simeq R_b,
  \qquad
  B\equiv k_{-}(R_b)=\bar A:\ x_B\simeq -R_b .
  \label{eq:app-packet-AB}
\end{equation}
With \(\mathcal K\equiv\mathcal C_{\rm disp}/\Delta R\), the four direct and mirror observables
associated with the pair are
\begin{align}
  \Iplus(R_b,t)
    &=\mathcal K\,[n_{A+}(t)-n_{A0}(t)],
  \label{eq:app-Iplus-direct}\\
  \Iminus(R_b,t)
    &=\mathcal K\,[n_{B0}(t)-n_{B-}(t)],
  \label{eq:app-Iminus-direct}\\
  \Iminus(-R_b,t)
    &=\mathcal K\,[n_{A0}(t)-n_{A-}(t)],
  \label{eq:app-Iminus-mirror}\\
  \Iplus(-R_b,t)
    &=\mathcal K\,[n_{B+}(t)-n_{B0}(t)].
  \label{eq:app-Iplus-mirror}
\end{align}
The signal at \(-R_b\) is not a second, independent reservoir. It is another measurement of a
population difference in the same packet that was modified at \(+R_b\). At \(R_b=0\), \(A=B\),
and one packet contributes both transitions at the same frequency.

\subsection{Elementary two-level equalization}

For two levels with a symmetric saturation process,
\begin{equation}
  J_{ab}=W_{ab}(n_a-n_b),\qquad \dot n_a=-J_{ab},\qquad \dot n_b=+J_{ab},\qquad
  \frac{\dd}{\dd t}(n_a-n_b)=-2W_{ab}(n_a-n_b).
  \label{eq:app-twolevel}
\end{equation}
The measured transition difference therefore decays at \(2W_{ab}\), twice the coefficient in the
symmetric current. The sign is automatic: if \(n_a<n_b\), population flows the other way.

\subsection{The RF currents with a finite homogeneous linewidth}

For a program with carrier commands \(U_j(t)\) and the RF gate \(u_{\rm RF}\in\{0,1\}\),
\begin{align}
  J_k^{+}
  &=u_{\rm RF}\,c_k\Bigl[\sum_jK_k(x_k-x_j)U_j\Bigr](n_{k+}-n_{k0}),
  \label{eq:app-Jplus}\\
  J_k^{-}
  &=u_{\rm RF}\,c_k\Bigl[\sum_jK_k(-x_k-x_j)U_j\Bigr](n_{k0}-n_{k-}),
  \label{eq:app-Jminus}
\end{align}
with \(K_k(\delta)=\kappa\Gamma_k/(\Gamma_k^2+\delta^2)\). The superscript on \(J^+_k\) denotes the
\(+\leftrightarrow0\) transition and on \(J^-_k\) the \(0\leftrightarrow-\) transition. In the
ideal-bin limit with one carrier at \(R_b\), only \(J^+_A\) and \(J^-_B\) are nonzero. In general,
every packet within a few \(\Gamma\) of any carrier, or of its mirror, is driven.

\subsection{The internal (exchange) currents}

With \(\tilde S_{ij}=\sqrt{f^{\rm ex}_if^{\rm ex}_j}\,S(x_i-x_j)\),
\(\tilde X_{ij}=\sqrt{f^{\rm ex}_if^{\rm ex}_j}\,S(x_i+x_j)(1-\tfrac12\delta_{ij})\), and
\(\tilde C_{ij}=\sqrt{f^{\rm ex}_if^{\rm ex}_j}\,C_{ij}\):
\begin{align}
G^{\rm SQ}_{k+}&=K_0\bigl[n_{k0}(\tilde Sn_+)_k-n_{k+}(\tilde Sn_0)_k\bigr],
\label{eq:app-sq-plus}\\
G^{\rm SQ}_{k-}&=-K_0\bigl[n_{k-}(\tilde Sn_0)_k-n_{k0}(\tilde Sn_-)_k\bigr],
\label{eq:app-sq-minus}\\
G^{\rm SQ}_{k0}&=-G^{\rm SQ}_{k+}-G^{\rm SQ}_{k-},
\\
G^{\rm X}_{k+}&=r_{\rm x}K_0\bigl[n_{k0}(\tilde Xn_0)_k-n_{k+}(\tilde Xn_-)_k\bigr],
\label{eq:app-x-plus}\\
G^{\rm X}_{k-}&=r_{\rm x}K_0\bigl[n_{k0}(\tilde Xn_0)_k-n_{k-}(\tilde Xn_+)_k\bigr],
\label{eq:app-x-minus}\\
G^{\rm X}_{k0}&=-G^{\rm X}_{k+}-G^{\rm X}_{k-},
\\
G^{\rm DQ}_{k+}&=-G^{\rm DQ}_{k-}=r_{\rm DQ}K_0\bigl[n_{k-}(\tilde Cn_+)_k-n_{k+}(\tilde Cn_-)_k\bigr],\qquad G^{\rm DQ}_{k0}=0 .
\label{eq:app-dq}
\end{align}
Each line follows from summing the pair currents of \cref{eq:Jsp,eq:Jsm,eq:Jx,eq:Jdq} over
partners \(j\). With DNP on, \(K_0\) carries the microwave factor. Summed over \(k\):
\begin{itemize}
\item \(\sum_kG^{\rm SQ}_{km}=0\) for each \(m\);
\item \(\sum_k(G^{\rm X}_{k+}-G^{\rm X}_{k-})=0\), using the symmetry of \(\tilde X\);
\item \(\sum_kG^{\rm DQ}_{k\pm}=0\).
\end{itemize}
These are the conservation laws of \cref{sec:stationary}, with no projection.

\subsection{Full population equations}

Collect all non-RF terms into
\begin{equation}
  G_{km}
  =G_{km}^{\rm SQ}
  +G_{km}^{\rm X}
  +G_{km}^{\rm DQ}
  +G_{km}^{\rm DNP}
  +G_{km}^{T_1},
  \label{eq:app-Gtotal}
\end{equation}
with
\begin{align}
  G_{km}^{\rm DNP}&=u_{\rm DNP}\,k_{\rm D}w^{\rm D}_k\bigl(\mu_kp^{\rm B}_m(P_{\rm sat})-n_{km}\bigr),\label{eq:app-dnp}\\
  G_{km}^{T_1}&=k_1f^{T_1}_k\bigl(\mu_kp^{\rm B}_m(P_{\rm eq})-n_{km}\bigr).\label{eq:app-T1}
\end{align}
The complete population equations are
\begin{align}
  \dot n_{k+}
  &=-J_k^{+}+G_{k+},
  \label{eq:app-full-nplus}\\
  \dot n_{k0}
  &=+J_k^{+}-J_k^{-}+G_{k0},
  \label{eq:app-full-nzero}\\
  \dot n_{k-}
  &=+J_k^{-}+G_{k-}.
  \label{eq:app-full-nminus}
\end{align}
The RF part conserves the packet spin count, because the three RF contributions sum to zero.

In the ideal-bin limit at \(R_b\neq0\), the directly addressed packets satisfy
\begin{align}
  \dot n_{A+}&=-J_A^{+}+G_{A+},
  &\dot n_{B+}&=G_{B+},
  \label{eq:app-six-1}\\
  \dot n_{A0}&=+J_A^{+}+G_{A0},
  &\dot n_{B0}&=-J_B^{-}+G_{B0},
  \label{eq:app-six-2}\\
  \dot n_{A-}&=G_{A-},
  &\dot n_{B-}&=+J_B^{-}+G_{B-}.
  \label{eq:app-six-3}
\end{align}
This is the simplest complete answer to the statement that ``the populations at \(R\) and \(-R\)
are different.'' With a finite linewidth, the same structure holds for every driven packet.

\subsection{Vector and tensor polarization at the population level}

\begin{equation}
  P(t)=\sum_k[n_{k+}-n_{k-}],\qquad
  Q(t)=\sum_k[n_{k+}-2n_{k0}+n_{k-}],
\end{equation}
\begin{align}
  \dot P&=-\sum_k(J_k^{+}+J_k^{-})
  +\sum_k\bigl(G^{\rm DNP}_{k+}-G^{\rm DNP}_{k-}+G^{T_1}_{k+}-G^{T_1}_{k-}\bigr),
  \label{eq:app-Pdot}\\
  \dot Q&=3\sum_k(J_k^{-}-J_k^{+})
  -3\sum_kG^{\rm X}_{k0}
  +\sum_k\bigl(G^{\rm DNP}_{k+}-2G^{\rm DNP}_{k0}+G^{\rm DNP}_{k-}\bigr)+(T_1\ \text{analog}).
  \label{eq:app-Qdot}
\end{align}
Exchange does not appear in \(\dot P\). Only the cross-branch channel appears in \(\dot Q\), where
\(-3\sum_kG^{\rm X}_{k0}=6\sum_{k,j}J^{({\rm x})}_{kj}\), which vanishes when \(p_0^2=p_+p_-\).

\subsection{Linearized local refill}

For a small \(\Iplus\) hole in packet \(k\) (\(p_{k+}=\bar p_+-\delta\), \(p_{k0}=\bar p_0+\delta\)),
with all partners at the Boltzmann fractions \(\bar p\), \cref{eq:app-sq-plus,eq:app-sq-minus}
give
\begin{align}
\dot p_{k+}\big|_{\rm SQ}&=K_0\bigl[(\tilde S\mu)_k-\tilde S_{kk}\mu_k\bigr](\bar p_++\bar p_0)\,\delta,\\
\dot p_{k0}\big|_{\rm SQ}&=-K_0\bigl[(\tilde S\mu)_k-\tilde S_{kk}\mu_k\bigr](\bar p_++\bar p_0+\bar p_-)\,\delta,
\end{align}
so that \(d_{k,+}=p_{k+}-p_{k0}\) relaxes at
\(\kappa^{\rm SQ}_k=K_0[(\tilde S\mu)_k-\tilde S_{kk}\mu_k](1-\bar p_-/2)\), which is
\cref{eq:kappaSQ}. The cross-branch and DQ contributions follow the same way, with \(\tilde X\) and
\(\tilde C\).

\subsection{From population equations to the plotted intensity rates}

For any packet,
\begin{align}
  \boxed{
  \dot\Iplus(x_k)
  =\mathcal K\left[
    -2J_k^{+}+J_k^{-}+G_{k+}-G_{k0}
  \right]}
  \label{eq:app-general-Idot-plus}\\[0.4em]
  \boxed{
  \dot\Iminus(-x_k)
  =\mathcal K\left[
    +J_k^{+}-2J_k^{-}+G_{k0}-G_{k-}
  \right].}
  \label{eq:app-general-Idot-minus}
\end{align}
At the monitor marker \(M=R_b\) (packets \(A\), \(B\)), the solid traces obey
\begin{align}
  \boxed{
  \frac{\dd\Iplus(R_b,t)}{\dd t}
  =\mathcal K\left[
    -2J_A^{+}+J_A^{-}+G_{A+}-G_{A0}
  \right]}
  \label{eq:app-direct-front-plus}\\[0.4em]
  \boxed{
  \frac{\dd\Iminus(R_b,t)}{\dd t}
  =\mathcal K\left[
    J_B^{+}-2J_B^{-}+G_{B0}-G_{B-}
  \right]}
  \label{eq:app-direct-front-minus}
\end{align}
and the dotted mirror traces obey
\begin{align}
  \boxed{
  \frac{\dd\Iminus(-R_b,t)}{\dd t}
  =\mathcal K\left[
    J_A^{+}-2J_A^{-}+G_{A0}-G_{A-}
  \right]}
  \label{eq:app-mirror-front-minus}\\[0.4em]
  \boxed{
  \frac{\dd\Iplus(-R_b,t)}{\dd t}
  =\mathcal K\left[
    -2J_B^{+}+J_B^{-}+G_{B+}-G_{B0}
  \right].}
  \label{eq:app-mirror-front-plus}
\end{align}
In the ideal-bin limit \(J^-_A=J^+_B=0\), and the first-release forms are recovered.

\subsection{RF-only limit and the one-half mirror-gain rule}

With all \(G\) set to zero and the ideal-bin limit,
\begin{align}
  \dot\Iplus(R_b)\big|_{\rm RF}
  &=-2\mathcal KJ_A^{+},
  &\dot\Iminus(-R_b)\big|_{\rm RF}
  &=+\mathcal KJ_A^{+},
  \label{eq:app-half-plus}\\
  \dot\Iminus(R_b)\big|_{\rm RF}
  &=-2\mathcal KJ_B^{-},
  &\dot\Iplus(-R_b)\big|_{\rm RF}
  &=+\mathcal KJ_B^{-},
  \label{eq:app-half-minus}
\end{align}
so the opposing-branch area gained is one half of the directly burned area lost
\cite{ClementKeller2023}. With a finite linewidth, \(J^-_A\) and \(J^+_B\) are nonzero near
\(R=0\), and the wings drive neighbouring packets. The exact relations are then
\cref{eq:app-direct-front-plus,eq:app-direct-front-minus,eq:app-mirror-front-minus,eq:app-mirror-front-plus}.

\subsection{Mechanism-by-mechanism intensity decomposition}
\label{app:breakdown}

\code{derivative(breakdown=True)} returns, for the monitored packets, the contribution of each
term to \(\dot\Iplus\), \(\dot\Iminus\) (direct and mirror), \(\dot P\) and \(\dot Q\):
\begin{align}
  \dot\Iplus_k
  &=\dot\Iplus_k^{\rm RF}
  +\dot\Iplus_k^{\rm SQ+}
  +\dot\Iplus_k^{\rm SQ-}
  +\dot\Iplus_k^{\rm X}
  +\dot\Iplus_k^{\rm DQ}
  +\dot\Iplus_k^{\rm DNP}
  +\dot\Iplus_k^{T_1},
\end{align}
with, for each term \(Y\), \(\dot\Iplus_k^{Y}=\mathcal K(G^Y_{k+}-G^Y_{k0})\) and
\(\dot\Iminus_k^{Y}=\mathcal K(G^Y_{k0}-G^Y_{k-})\). These are the natural quantities to compare
with measured short-time slopes. They are also the tool for separating RF response from
replenishment in the peak--pedestal tests.

\subsection{Numerical time step and constraints}

The equations are advanced with the adaptive, event-split Euler step of \cref{eq:euler}. A
population below \(-10^{-13}\mu_k\) is an error. Tiny rounding negatives are floored at
\(10^{-30}\), and each packet is rescaled to \(\mu_k\):
\begin{equation}
  \bm n_k
  \leftarrow
  \mu_k
  \frac{\bm n_k}{\sum_m n_{km}}.
\end{equation}
These operations are numerical safeguards, not physical source terms. The step bound
\(h\le0.2/\lambda_{\max}\) guarantees positivity for the linear terms, independently of the
command amplitude, so commands are never silently corrected by the clamp.

\subsection{Operational summary}

At every substep, the software performs the following sequence:
\begin{enumerate}
  \item determine the next pulse edge and the active command vector \(U\);
  \item compute \(W^\pm_k\) from \(U\) by convolution, and the RF currents \(J^\pm_k\) from the current population differences;
  \item compute the exchange fields \(Sn\), \(Xn\) and \(Cn\) by FFT and \(O(N)\) sums, and the currents \(G^{\rm SQ}\), \(G^{\rm X}\), \(G^{\rm DQ}\);
  \item add the DNP and \(T_1\) terms (and, in the extended engine, the reservoir terms and \(\dot\beta\));
  \item choose \(h\) from the rate bound and the next edge, advance, check positivity, renormalize, and accumulate exposure;
  \item project the populations to \(\Iplus(R,t)\) and \(\Iminus(R,t)\) through \cref{eq:app-Iplus-direct,eq:app-Iminus-direct,eq:app-Iminus-mirror,eq:app-Iplus-mirror} and their full-grid equivalents.
\end{enumerate}

%==================================================================================================
\section{Interpretation of the main rate controls}

\begin{longtable}{>{\ttfamily}L{0.33\textwidth} L{0.08\textwidth} L{0.51\textwidth}}
\toprule
Parameter & Symbol & Recommended interpretation\\
\midrule
rf\_spectral\_strength & \(\kappa\) & RF rate scale per unit command, in \(R\) units; \(\kappa/\Gamma\) is the on-resonance rate per \(U\).\\
rf\_homogeneous\_hwhm\_R & \(\Gamma\) & Homogeneous (\(T_2\)) half-width of a packet's transitions; sets hole width.\\
rf\_peak\_selectivity & \(s\) & Density exponent of per-spin RF coupling inside the peak bands; not fitted.\\
rf\_peak\_center\_R,\newline rf\_peak\_half\_width\_R,\newline rf\_transition\_width\_R & \(R_c\) & Location and taper of the peak bands.\\
rf\_coupling\_pedestal\_gain & \(G_{\rm p}\) & Per-spin coupling outside the peak bands.\\
X\_density\_exponent,\newline X\_theta\_p2/p4,\newline X\_od\_factor & \(e_X\),\newline \(a_{2,4}^X\),\newline \(m^X_{\rm OD}\) & Mechanism profile, \cref{eq:profile}, for X in rf\_coupling, rf\_width, exchange, dnp, t1.\\
diffusion\_scale & \(K_0\) & Flip-flop rate scale; refill rate is \(K_0\) times the partner count.\\
zq\_width\_R, diffusion\_overlap & \(w\), \(S\) & ZQ overlap width and shape; narrow \(w\) makes transport local and density-dependent.\\
cross\_branch\_ratio & \(r_{\rm x}\) & Cross-branch / same-branch; 1 from equal matrix elements; the only channel that relaxes \(Q\) at fixed \(P\).\\
double\_quantum\_ratio & \(r_{\rm DQ}\) & Effective second-order DQ exchange; nonlocal vector transport.\\
microwave\_diffusion\_factor & --- & Multiplies \(K_0\) while DNP is on.\\
orientation\_corr\_fraction, \_width\_deg & \(f_\theta,\sigma_\theta\) & Optional orientation--space correlation of exchange partners.\\
kernel\_cutoff\_widths & --- & Optional hard cutoff of the overlap in units of \(w\); 0 keeps all tails.\\
dnp\_rate, p\_dnp\_sat & \(k_{\rm D},P_{\rm sat}\) & DNP build rate and target.\\
capacity\_rate\_power, \_clip & \(\alpha_{\rm D},c\) & DNP density weighting (legacy RF modes also use it).\\
t1\_rate, t1\_p\_eq & \(k_1,P_{\rm eq}\) & Spin-lattice relaxation and its target.\\
site\_resolved, exchange\_intersite\_factor & \(M_{12}\) & Separate C--D/O--D populations and their mutual exchange.\\
dipolar\_enabled, dipolar\_capacity\_ratio & \(c_d\) & Dipolar reservoir and its heat capacity.\\
dipolar\_exchange\_coupling & \(\xi\) & Flip-flop coupling to the reservoir.\\
dipolar\_relaxation\_rate, dipolar\_beta\_lattice & \(r_{1d},\beta_L\) & Reservoir relaxation.\\
afp\_efficiency & \(\varepsilon\) & AFP swap efficiency.\\
dt & \(\delta t\) & Material step; substeps are chosen automatically.\\
\bottomrule
\end{longtable}

%==================================================================================================
\begin{thebibliography}{99}

\bibitem[Pake(1948)]{Pake1948}
G.~E. Pake,
``Nuclear Resonance Absorption in Hydrated Crystals: Fine Structure of the Proton Line,''
\emph{Journal of Chemical Physics} \textbf{16}(4), 327--336 (1948),
\href{https://doi.org/10.1063/1.1746878}{doi:10.1063/1.1746878}.

\bibitem[Bloch(1946)]{Bloch1946}
F. Bloch,
``Nuclear Induction,''
\emph{Physical Review} \textbf{70}(7--8), 460--474 (1946),
\href{https://doi.org/10.1103/PhysRev.70.460}{doi:10.1103/PhysRev.70.460}.

\bibitem[Abragam(1961)]{Abragam1961}
A. Abragam,
\emph{The Principles of Nuclear Magnetism},
Oxford University Press, Oxford (1961).

\bibitem[Slichter(1990)]{Slichter1990}
C.~P. Slichter,
\emph{Principles of Magnetic Resonance}, 3rd ed.,
Springer, Berlin (1990).

\bibitem[Wasylishen et al.(2012)]{Wasylishen2012}
R.~E. Wasylishen, S.~E. Ashbrook, and S. Wimperis (eds.),
\emph{NMR of Quadrupolar Nuclei in Solid Materials},
John Wiley \& Sons, Chichester (2012).

\bibitem[Cohen and Reif(1957)]{CohenReif1957}
M.~H. Cohen and F. Reif,
``Quadrupole Effects in Nuclear Magnetic Resonance Studies of Solids,''
\emph{Solid State Physics} \textbf{5}, 321--438 (1957).

\bibitem[Redfield(1955)]{Redfield1955}
A.~G. Redfield,
``Nuclear Magnetic Resonance Saturation and Rotary Saturation in Solids,''
\emph{Physical Review} \textbf{98}(6), 1787--1809 (1955).

\bibitem[Provotorov(1962)]{Provotorov1962}
B.~N. Provotorov,
``Magnetic Resonance Saturation in Crystals,''
\emph{Soviet Physics JETP} \textbf{14}, 1126 (1962).

\bibitem[Goldman(1970)]{Goldman1970}
M. Goldman,
\emph{Spin Temperature and Nuclear Magnetic Resonance in Solids},
Clarendon Press, Oxford (1970).

\bibitem[Abragam and Goldman(1978)]{AbragamGoldman1978}
A. Abragam and M. Goldman,
``Principles of Dynamic Nuclear Polarisation,''
\emph{Reports on Progress in Physics} \textbf{41}(3), 395--467 (1978),
\href{https://doi.org/10.1088/0034-4885/41/3/002}{doi:10.1088/0034-4885/41/3/002}.

\bibitem[Abragam and Goldman(1982)]{AbragamGoldman1982}
A. Abragam and M. Goldman,
\emph{Nuclear Magnetism: Order and Disorder},
Clarendon Press, Oxford (1982).

\bibitem[Borghini(1968)]{Borghini1968}
M. Borghini,
``Spin-Temperature Model of Nuclear Dynamic Polarization Using Free Radicals,''
\emph{Physical Review Letters} \textbf{20}, 419 (1968).

\bibitem[Wenckebach(2016)]{Wenckebach2016}
W.~T. Wenckebach,
\emph{Essentials of Dynamic Nuclear Polarization},
Spindrift Publications, The Netherlands (2016).

\bibitem[Khutsishvili(1962)]{Khutsishvili1962}
G.~R. Khutsishvili,
``Spin Diffusion and Magnetic Relaxation of Nuclei,''
\emph{Soviet Physics JETP} \textbf{15}(5), 909--913 (1962),
\href{https://jetp.ras.ru/cgi-bin/dn/e_015_05_0909.pdf}{publisher PDF}.

\bibitem[Khutsishvili(1966)]{Khutsishvili1966}
G.~R. Khutsishvili,
``Spin Diffusion,''
\emph{Soviet Physics Uspekhi} \textbf{8}, 743 (1966).

\bibitem[Suter and Ernst(1985)]{SuterErnst1985}
D. Suter and R.~R. Ernst,
``Spin Diffusion in Resolved Solid-State NMR Spectra,''
\emph{Physical Review B} \textbf{32}(9), 5608--5627 (1985),
\href{https://doi.org/10.1103/PhysRevB.32.5608}{doi:10.1103/PhysRevB.32.5608}.

\bibitem[Eckman(1982)]{Eckman1982}
R.~R. Eckman,
\emph{Hydrogen and Deuterium NMR of Solids by Magic-Angle Spinning},
Ph.D. thesis, University of California, Berkeley;
Lawrence Berkeley Laboratory report LBL-14200 (1982),
Sec.~III.G, pp.~133--146,
\href{https://doi.org/10.2172/6645809}{doi:10.2172/6645809}.

\bibitem[Walder et al.(2019)]{Walder2019}
B.~J. Walder, N.~A. Prisco, F.~M. Paruzzo, J.~R. Yarava,
B.~F. Chmelka, and L. Emsley,
``Measurement of Proton Spin Diffusivity in Hydrated Cementitious Solids,''
\emph{Journal of Physical Chemistry Letters} \textbf{10}(17),
5064--5069 (2019),
\href{https://doi.org/10.1021/acs.jpclett.9b01861}{doi:10.1021/acs.jpclett.9b01861}.

\bibitem[Garwood and DelaBarre(2001)]{GarwoodDelaBarre2001}
M. Garwood and L. DelaBarre,
``The Return of the Frequency Sweep: Designing Adiabatic Pulses for Contemporary NMR,''
\emph{Journal of Magnetic Resonance} \textbf{153}, 155--177 (2001).

\bibitem[Dulya et al.(1997)]{Dulya1997}
C. Dulya \emph{et al.},
``A Line-Shape Analysis for Spin-1 NMR Signals,''
\emph{Nuclear Instruments and Methods in Physics Research Section A} \textbf{398}, 109--125 (1997).

\bibitem[Keller(2017)]{Keller2017}
D. Keller,
``Modeling Alignment Enhancement for Solid Polarized Targets,''
\emph{European Physical Journal A} \textbf{53}, 155 (2017),
\href{https://doi.org/10.1140/epja/i2017-12344-0}{doi:10.1140/epja/i2017-12344-0}.

\bibitem[Keller et al.(2020)]{KellerCrabbDay2020}
D. Keller, D. Crabb, and D. Day,
``Enhanced Tensor Polarization in Solid-State Targets,''
\emph{Nuclear Instruments and Methods in Physics Research Section A} \textbf{981}, 164504 (2020),
\href{https://doi.org/10.1016/j.nima.2020.164504}{doi:10.1016/j.nima.2020.164504}.

\bibitem[Clement and Keller(2023)]{ClementKeller2023}
J. Clement and D. Keller,
``Manipulation of Spin-1 Solid-State Targets,''
\emph{Nuclear Instruments and Methods in Physics Research Section A} \textbf{1050}, 168177 (2023),
\href{https://doi.org/10.1016/j.nima.2023.168177}{doi:10.1016/j.nima.2023.168177}.

\bibitem[Keller(2026)]{Keller2026}
D. Keller,
``Selective Semi-Saturation Measurement of Spin-1,''
\emph{European Physical Journal A} \textbf{62}, 36 (2026),
\href{https://doi.org/10.1140/epja/s10050-026-01790-y}{doi:10.1140/epja/s10050-026-01790-y}.

\bibitem[Giles and Pierce(2000)]{GilesPierce2000}
M.~B. Giles and N.~A. Pierce,
``An Introduction to the Adjoint Approach to Design,''
\emph{Flow, Turbulence and Combustion} \textbf{65}, 393--415 (2000).

\bibitem[Plessix(2006)]{Plessix2006}
R.-E. Plessix,
``A Review of the Adjoint-State Method for Computing the Gradient of a Functional with Geophysical Applications,''
\emph{Geophysical Journal International} \textbf{167}, 495--503 (2006).

\bibitem[Byrd et al.(1995)]{Byrd1995}
R.~H. Byrd, P. Lu, J. Nocedal, and C. Zhu,
``A Limited Memory Algorithm for Bound Constrained Optimization,''
\emph{SIAM Journal on Scientific Computing} \textbf{16}, 1190--1208 (1995).

\bibitem[Kraft(1988)]{Kraft1988}
D. Kraft,
``A Software Package for Sequential Quadratic Programming,''
DFVLR-FB 88-28, Deutsche Forschungs- und Versuchsanstalt f\"ur Luft- und Raumfahrt (1988).

\end{thebibliography}

\end{document}
